% !TeX program = lualatex
% ALIUS Bulletin native visual reconstruction source.
% Generated from off-repo reference layout data; does not include or input a pre-existing article PDF.
\ifdefined\ALIUSIssueBuild
\else
\documentclass{article}
\usepackage[paperwidth=595.2756bp,paperheight=841.8898bp,margin=0bp]{geometry}
\usepackage{graphicx}
\usepackage{xcolor}
\usepackage{tikz}
\usepackage{iftex}
\usepackage[hidelinks]{hyperref}
\ifPDFTeX
  \usepackage[T1]{fontenc}
  \usepackage[utf8]{inputenc}
  \usepackage{textcomp}
  \usepackage{amssymb}
\else
  \usepackage{fontspec}
\fi
\pagestyle{empty}
\fi
\ifPDFTeX
  \providecommand{\ALIUSDeclareUnicodeFallbacks}{%
    \DeclareUnicodeCharacter{00AC}{\ensuremath{\neg}}%
    \DeclareUnicodeCharacter{00BE}{\ensuremath{\frac{3}{4}}}%
    \DeclareUnicodeCharacter{0101}{\=a}%
    \DeclareUnicodeCharacter{0107}{\'c}%
    \DeclareUnicodeCharacter{012B}{\=\i}%
    \DeclareUnicodeCharacter{0131}{\i}%
    \DeclareUnicodeCharacter{0142}{\l}%
    \DeclareUnicodeCharacter{015B}{\'s}%
    \DeclareUnicodeCharacter{016B}{\=u}%
    \DeclareUnicodeCharacter{025C}{q}%
    \DeclareUnicodeCharacter{0278}{t}%
    \DeclareUnicodeCharacter{02AE}{Th}%
    \DeclareUnicodeCharacter{02B0}{ff}%
    \DeclareUnicodeCharacter{02B4}{ffi}%
    \DeclareUnicodeCharacter{02BE}{ft}%
    \DeclareUnicodeCharacter{02BF}{fi}%
    \DeclareUnicodeCharacter{02C0}{fl}%
    \DeclareUnicodeCharacter{02C1}{fi}%
    \DeclareUnicodeCharacter{02D9}{Th}%
    \DeclareUnicodeCharacter{02EF}{gy}%
    \DeclareUnicodeCharacter{0394}{\ensuremath{\Delta}}%
    \DeclareUnicodeCharacter{03B2}{\ensuremath{\beta}}%
    \DeclareUnicodeCharacter{03BA}{\ensuremath{\kappa}}%
    \DeclareUnicodeCharacter{068D}{-}%
    \DeclareUnicodeCharacter{08BC}{ti}%
    \DeclareUnicodeCharacter{099E}{ti}%
    \DeclareUnicodeCharacter{1E43}{\d{m}}%
    \DeclareUnicodeCharacter{1E47}{\d{n}}%
    \DeclareUnicodeCharacter{1E63}{\d{s}}%
    \DeclareUnicodeCharacter{201C}{``}%
    \DeclareUnicodeCharacter{201D}{''}%
    \DeclareUnicodeCharacter{2018}{`}%
    \DeclareUnicodeCharacter{2019}{'}%
    \DeclareUnicodeCharacter{2260}{\ensuremath{\neq}}%
    \DeclareUnicodeCharacter{25A1}{\ensuremath{\square}}%
  }%
  \ifdefined\ALIUSIssueBuild\else\ALIUSDeclareUnicodeFallbacks\fi
  \providecommand{\ALIUSFontLatoLight}{\fontfamily{lato-LF}\fontseries{l}\selectfont}
  \providecommand{\ALIUSFontLatoRegular}{\fontfamily{lato-LF}\fontseries{m}\selectfont}
  \providecommand{\ALIUSFontLatoItalic}{\fontfamily{lato-LF}\fontseries{m}\itshape}
  \providecommand{\ALIUSFontLatoLightItalic}{\fontfamily{lato-LF}\fontseries{l}\itshape}
  \providecommand{\ALIUSFontCormorantRegular}{\fontfamily{CormorantGaramond-LF}\fontseries{m}\selectfont}
  \providecommand{\ALIUSFontCormorantLight}{\fontfamily{CormorantGaramond-LF}\fontseries{l}\selectfont}
  \providecommand{\ALIUSFontCormorantMedium}{\fontfamily{CormorantGaramond-LF}\fontseries{medium}\selectfont}
  \providecommand{\ALIUSFontCormorantBold}{\fontfamily{CormorantGaramond-LF}\fontseries{b}\selectfont}
  \providecommand{\ALIUSFontCormorantItalic}{\fontfamily{CormorantGaramond-LF}\fontseries{m}\itshape}
  \providecommand{\ALIUSFontTimes}{\fontfamily{ptm}\selectfont}
  \providecommand{\ALIUSFontCalibri}{\fontfamily{phv}\selectfont}
  \providecommand{\ALIUSFontCambria}{\fontfamily{ptm}\selectfont}
  \providecommand{\ALIUSFontSymbolFallback}{\fontfamily{ptm}\selectfont}
\else
  \providecommand{\ALIUSFontLatoLight}{\fontspec{Lato Light}}
  \providecommand{\ALIUSFontLatoRegular}{\fontspec{Lato Regular}}
  \providecommand{\ALIUSFontLatoItalic}{\fontspec{Lato Italic}}
  \providecommand{\ALIUSFontLatoLightItalic}{\fontspec{Lato Light Italic}}
  \providecommand{\ALIUSFontCormorantRegular}{\fontspec{Cormorant Garamond Regular}}
  \providecommand{\ALIUSFontCormorantLight}{\fontspec{Cormorant Garamond Light}}
  \providecommand{\ALIUSFontCormorantMedium}{\fontspec{Cormorant Garamond Medium}}
  \providecommand{\ALIUSFontCormorantBold}{\fontspec{Cormorant Garamond Bold}}
  \providecommand{\ALIUSFontCormorantItalic}{\fontspec{Cormorant Garamond Italic}}
  \providecommand{\ALIUSFontTimes}{\fontspec{Times New Roman}}
  \providecommand{\ALIUSFontCalibri}{\fontspec{Calibri}}
  \providecommand{\ALIUSFontCambria}{\fontspec{Cambria}}
  \providecommand{\ALIUSFontSymbolFallback}{\fontspec{Times New Roman}}
\fi
\providecommand{\ALIUSPullQuoteOpen}{\textquotedblleft}
\providecommand{\ALIUSPullQuoteClose}{\textquotedblright}
\providecommand{\ALIUSPullQuoteBlockAt}[4]{%
  \node[anchor=north west,inner sep=0pt,outer sep=0pt,xshift=12bp,text=ALIUSC595959,font=\ALIUSFontLatoLight\fontsize{15.000bp}{18.000bp}\selectfont\hyphenpenalty10000\exhyphenpenalty10000\relax] (ALIUSPullQuoteBody) at (#1bp,-#2bp) {\begin{tabular}{@{}c@{}}#4\end{tabular}};%
  \node[anchor=north east,inner sep=0pt,outer sep=0pt,text=ALIUSC7F7F7F,xshift=-6bp] at (ALIUSPullQuoteBody.north west) {%
    {\ALIUSFontCormorantLight\fontsize{32.000bp}{32.000bp}\selectfont\ALIUSPullQuoteOpen}%
  };%
  \node[anchor=south west,inner sep=0pt,outer sep=0pt,text=ALIUSC7F7F7F,xshift=6bp,yshift=-7bp] at (ALIUSPullQuoteBody.south east) {%
    {\ALIUSFontCormorantLight\fontsize{32.000bp}{32.000bp}\selectfont\ALIUSPullQuoteClose}%
  };%
}
\providecommand{\ALIUSNotableQuoteAt}[4]{%
  \ALIUSPullQuoteBlockAt{#1}{#2}{#3}{#4}%
}
\providecommand{\ALIUSMaybeNotableQuoteAt}[4]{%
  \if\relax\detokenize{#4}\relax\else\ALIUSNotableQuoteAt{#1}{#2}{#3}{#4}\fi%
}
\providecommand{\ALIUSRefAnchor}[1]{\hypertarget{#1}{}}
\providecommand{\ALIUSCitationLink}[2]{\hyperlink{#1}{#2}}
\definecolor{ALIUSC000000}{HTML}{000000}
\definecolor{ALIUSC0000FF}{HTML}{0000FF}
\definecolor{ALIUSC1F8135}{HTML}{1F8135}
\definecolor{ALIUSC595959}{HTML}{595959}
\definecolor{ALIUSC767171}{HTML}{767171}
\definecolor{ALIUSCBFBFBF}{HTML}{BFBFBF}
\definecolor{ALIUSCFF0000}{HTML}{FF0000}
\definecolor{ALIUSCFFFFFF}{HTML}{FFFFFF}
\providecommand{\ALIUSPlacedText}[6]{%
  \node[anchor=north west,inner sep=0pt,outer sep=0pt,text=#4] at (#1bp,-#2bp) {%
    \resizebox{#3bp}{!}{{#5\fontsize{#6bp}{#6bp}\selectfont #6\kern0pt}}%
  };%
}
\providecommand{\ALIUSPlacedTextContent}[7]{%
  \node[anchor=north west,inner sep=0pt,outer sep=0pt,text=#4] at (#1bp,-#2bp) {%
    \resizebox{#3bp}{!}{{#5\fontsize{#6bp}{#6bp}\selectfont #7}}%
  };%
}
\ifdefined\ALIUSIssueBuild
\else
\begin{document}
\fi
% Page 1
\thispagestyle{empty}
\noindent\begin{tikzpicture}[remember picture,overlay,x=1bp,y=1bp,shift={(current page.north west)}]
  \fill[ALIUSC1F8135] (401.031bp,-224.391bp) rectangle ++(65.520bp,-0.240bp);
  \fill[ALIUSC1F8135] (401.031bp,-284.391bp) rectangle ++(114.000bp,-0.240bp);
  \fill[ALIUSCBFBFBF] (395.511bp,-201.351bp) rectangle ++(0.480bp,-69.600bp);
  \fill[ALIUSC0000FF] (194.151bp,-320.151bp) rectangle ++(162.720bp,-0.480bp);
  \fill[ALIUSCBFBFBF] (395.511bp,-270.951bp) rectangle ++(0.480bp,-2.160bp);
  \fill[ALIUSCBFBFBF] (395.511bp,-273.111bp) rectangle ++(0.480bp,-48.240bp);
  \fill[ALIUSCFFFFFF] (90.951bp,-395.031bp) rectangle ++(414.960bp,-14.400bp);
  \fill[ALIUSCFFFFFF] (90.951bp,-409.431bp) rectangle ++(414.960bp,-14.400bp);
  \fill[ALIUSCFFFFFF] (105.111bp,-423.831bp) rectangle ++(400.800bp,-14.400bp);
  \fill[ALIUSCFFFFFF] (90.951bp,-438.231bp) rectangle ++(414.960bp,-14.400bp);
  \fill[ALIUSCFFFFFF] (90.951bp,-452.631bp) rectangle ++(414.960bp,-14.400bp);
  \fill[ALIUSCFFFFFF] (90.951bp,-467.031bp) rectangle ++(414.960bp,-14.400bp);
  \fill[ALIUSCFFFFFF] (90.951bp,-481.431bp) rectangle ++(414.960bp,-14.400bp);
  \fill[ALIUSCFFFFFF] (90.951bp,-495.831bp) rectangle ++(215.520bp,-15.120bp);
  \fill[ALIUSCFFFFFF] (90.951bp,-566.151bp) rectangle ++(414.960bp,-15.120bp);
  \fill[ALIUSCFFFFFF] (90.951bp,-583.431bp) rectangle ++(414.960bp,-15.120bp);
  \fill[ALIUSCFFFFFF] (90.951bp,-600.711bp) rectangle ++(414.960bp,-15.120bp);
  \fill[ALIUSCFFFFFF] (90.951bp,-617.751bp) rectangle ++(414.960bp,-15.120bp);
  \fill[ALIUSCFFFFFF] (90.951bp,-635.031bp) rectangle ++(414.960bp,-15.120bp);
  \fill[ALIUSCFFFFFF] (90.951bp,-652.311bp) rectangle ++(268.080bp,-15.120bp);
  \fill[ALIUSCFFFFFF] (90.951bp,-686.871bp) rectangle ++(414.960bp,-17.040bp);
  \fill[ALIUSCFFFFFF] (90.951bp,-706.311bp) rectangle ++(414.960bp,-17.040bp);
  \fill[ALIUSCFFFFFF] (90.951bp,-725.751bp) rectangle ++(414.960bp,-17.280bp);
  \fill[ALIUSCFFFFFF] (90.951bp,-745.431bp) rectangle ++(414.960bp,-17.040bp);
  \ALIUSPlacedTextContent{90.951}{72.970}{358.794}{ALIUSC000000}{\ALIUSFontLatoRegular}{19.920}{Internally Generated Conscious Activity:}
  \ALIUSPlacedTextContent{90.951}{108.970}{346.427}{ALIUSC000000}{\ALIUSFontLatoLight}{19.920}{Reflections upon (lucid) dreaming, mind-}
  \ALIUSPlacedTextContent{90.951}{132.970}{230.116}{ALIUSC000000}{\ALIUSFontLatoLight}{19.920}{wandering and meditation}
  \ALIUSPlacedTextContent{401.031}{201.254}{76.013}{ALIUSC000000}{\ALIUSFontLatoRegular}{11.040}{Benjamin Baird}
  \ALIUSPlacedTextContent{88.791}{201.320}{130.764}{ALIUSC000000}{\ALIUSFontLatoLight}{16.080}{An interview with}
  \ALIUSPlacedTextContent{401.031}{214.429}{72.891}{ALIUSC1F8135}{\ALIUSFontLatoLight}{9.120}{bbaird@wisc.edu}
  \ALIUSPlacedTextContent{88.791}{220.557}{128.693}{ALIUSC000000}{\ALIUSFontLatoLight}{18.960}{Benjamin Baird}
  \ALIUSPlacedTextContent{401.031}{225.229}{91.376}{ALIUSC000000}{\ALIUSFontLatoLight}{9.120}{University of Madison-}
  \ALIUSPlacedTextContent{401.031}{236.029}{64.125}{ALIUSC000000}{\ALIUSFontLatoLight}{9.120}{Wisconsin, USA}
  \ALIUSPlacedTextContent{88.791}{255.279}{119.346}{ALIUSC000000}{\ALIUSFontLatoLight}{12.960}{by Matthieu Koroma}
  \ALIUSPlacedTextContent{401.031}{261.254}{88.206}{ALIUSC000000}{\ALIUSFontLatoRegular}{11.040}{Matthieu Koroma}
  % ALIUS normalized citation panel begin
  % ALIUS normalized citation panel text: Cite as: Baird, B. \& Koroma, M. (2020). Internally generated conscious activity: Reflections upon (lucid) dreaming, mind-wandering and meditation. An interview with Benjamin Baird by Matthieu Koroma. ALIUS Bulletin, 4, 24-38,
  \fill[white] (87.291bp,-271.242bp) rectangle ++(305.240bp,-51.300bp);
  \node[anchor=north west,inner sep=0pt,outer sep=0pt,text=ALIUSC000000] at (88.791bp,-273.242bp) {\begin{minipage}{302.240bp}{\ALIUSFontLatoLight\fontsize{9.200bp}{10.700bp}\selectfont\raggedright Cite as: Baird, B. \& Koroma, M. (2020). Internally generated conscious activity: Reflections upon (lucid) dreaming, mind-wandering and meditation. An interview with Benjamin Baird by Matthieu Koroma. ALIUS Bulletin, 4, 24-38, \mbox{\textcolor{ALIUSC1F8135}{\href{https://doi.org/10.34700/h9bs-xs28}{https://doi.org/10.34700/h9bs-xs28}}}\par}\end{minipage}};
  % ALIUS normalized citation panel end
  \ALIUSPlacedTextContent{401.031}{274.429}{116.378}{ALIUSC1F8135}{\ALIUSFontLatoLight}{9.120}{mkoroma@ens-paris-saclay.fr}
  \ALIUSPlacedTextContent{401.031}{285.229}{105.901}{ALIUSC000000}{\ALIUSFontLatoLight}{9.120}{Ecole Normale Supérieure,}
  \ALIUSPlacedTextContent{401.031}{296.029}{51.966}{ALIUSC000000}{\ALIUSFontLatoLight}{9.120}{Paris, France}
  \ALIUSPlacedTextContent{90.951}{366.412}{55.719}{ALIUSC000000}{\ALIUSFontLatoRegular}{13.920}{Abstract}
  \ALIUSPlacedTextContent{90.951}{394.947}{418.269}{ALIUSC000000}{\ALIUSFontLatoLight}{12.000}{Certain conscious states such as dreaming reveal that conscious activity can be}
  \ALIUSPlacedTextContent{90.951}{409.347}{418.274}{ALIUSC000000}{\ALIUSFontLatoLight}{12.000}{to a large extent internally generated rather than being driven by sensory stimuli.}
  \ALIUSPlacedTextContent{90.951}{423.747}{418.275}{ALIUSC000000}{\ALIUSFontLatoLight}{12.000}{In this interview, psychologist and neuroscientist Benjamin Baird discusses the}
  \ALIUSPlacedTextContent{90.951}{438.147}{418.271}{ALIUSC000000}{\ALIUSFontLatoLight}{12.000}{developments of scientific research on these conscious phenomena including}
  \ALIUSPlacedTextContent{90.951}{452.547}{418.272}{ALIUSC000000}{\ALIUSFontLatoLight}{12.000}{dreaming, mind-wandering and meditation and how they interrelate. Lucid}
  \ALIUSPlacedTextContent{90.951}{466.947}{418.267}{ALIUSC000000}{\ALIUSFontLatoLight}{12.000}{dreaming, the ability to become aware that one is dreaming while in a dream, is}
  \ALIUSPlacedTextContent{90.951}{481.347}{418.264}{ALIUSC000000}{\ALIUSFontLatoLight}{12.000}{highlighted as a unique way to gain experimental control over internally}
  \ALIUSPlacedTextContent{90.951}{495.987}{212.874}{ALIUSC000000}{\ALIUSFontLatoLight}{12.000}{generated conscious activity during sleep}
  \ALIUSPlacedTextContent{303.813}{497.487}{2.628}{ALIUSC000000}{\ALIUSFontCormorantRegular}{12.000}{.}
  \ALIUSPlacedTextContent{90.951}{522.867}{47.340}{ALIUSC000000}{\ALIUSFontLatoItalic}{12.000}{keywords}
  \ALIUSPlacedTextContent{138.289}{522.867}{349.726}{ALIUSC000000}{\ALIUSFontLatoLightItalic}{12.000}{: dreaming, lucid dreaming, mind-wandering, meditation, metacognition}
  \ALIUSPlacedTextContent{90.951}{566.373}{418.407}{ALIUSC1F8135}{\ALIUSFontLatoLight}{12.357}{You hold a PhD in cognitive neuroscience and currently are a postdoctoral}
  \ALIUSPlacedTextContent{90.951}{583.653}{418.304}{ALIUSC1F8135}{\ALIUSFontLatoLight}{12.357}{research fellow at the Wisconsin Institute for Sleep and Consciousness in the}
  \ALIUSPlacedTextContent{90.951}{600.933}{418.256}{ALIUSC1F8135}{\ALIUSFontLatoLight}{12.357}{School of Medicine and Public Health at the University of Wisconsin, Madison.}
  \ALIUSPlacedTextContent{90.951}{617.973}{418.395}{ALIUSC1F8135}{\ALIUSFontLatoLight}{12.357}{In your research, you have investigated topics such as mind-wandering,}
  \ALIUSPlacedTextContent{90.951}{635.253}{418.374}{ALIUSC1F8135}{\ALIUSFontLatoLight}{12.357}{metacognition, meditation and dreaming. Could you introduce the scientific}
  \ALIUSPlacedTextContent{90.951}{652.533}{267.922}{ALIUSC1F8135}{\ALIUSFontLatoLight}{12.357}{questions that you are targeting in your research?}
  \ALIUSPlacedTextContent{90.951}{686.969}{418.185}{ALIUSC000000}{\ALIUSFontCormorantRegular}{13.920}{Since I was young I have been interested in conscious experience and how it}
  \ALIUSPlacedTextContent{90.951}{706.409}{418.235}{ALIUSC000000}{\ALIUSFontCormorantRegular}{13.920}{arises in nature. Overall, my long-term research interests extend to broad}
  \ALIUSPlacedTextContent{90.951}{725.849}{418.304}{ALIUSC000000}{\ALIUSFontCormorantRegular}{13.920}{scientific questions (which I grant I will likely never see answered in my}
  \ALIUSPlacedTextContent{90.951}{745.529}{418.299}{ALIUSC000000}{\ALIUSFontCormorantRegular}{13.920}{lifetime but I pursue them nonetheless!) including: how subjectivity arises in}
  \ALIUSPlacedTextContent{90.951}{793.143}{122.397}{ALIUSC767171}{\ALIUSFontCormorantLight}{12.000}{ALIUS Bulletin n°4 (2020)}
  \ALIUSPlacedTextContent{387.435}{793.143}{120.728}{ALIUSC767171}{\ALIUSFontCormorantLight}{12.000}{aliusresearch.org/bulletin}
\end{tikzpicture}
\null
\newpage
% Page 2
\thispagestyle{empty}
\noindent\begin{tikzpicture}[remember picture,overlay,x=1bp,y=1bp,shift={(current page.north west)}]
  \fill[ALIUSCFFFFFF] (90.951bp,-72.951bp) rectangle ++(414.960bp,-17.040bp);
  \fill[ALIUSCFFFFFF] (90.951bp,-92.391bp) rectangle ++(414.960bp,-17.040bp);
  \fill[ALIUSCFFFFFF] (90.951bp,-111.831bp) rectangle ++(414.960bp,-17.280bp);
  \fill[ALIUSCFFFFFF] (90.951bp,-131.511bp) rectangle ++(414.960bp,-17.040bp);
  \fill[ALIUSCFFFFFF] (90.951bp,-150.951bp) rectangle ++(414.960bp,-17.040bp);
  \fill[ALIUSCFFFFFF] (90.951bp,-170.391bp) rectangle ++(414.960bp,-17.040bp);
  \fill[ALIUSCFFFFFF] (90.951bp,-189.831bp) rectangle ++(414.960bp,-17.280bp);
  \fill[ALIUSCFFFFFF] (90.951bp,-209.511bp) rectangle ++(414.960bp,-17.040bp);
  \fill[ALIUSCFFFFFF] (90.951bp,-228.951bp) rectangle ++(414.960bp,-17.040bp);
  \fill[ALIUSCFFFFFF] (90.951bp,-248.391bp) rectangle ++(414.960bp,-17.040bp);
  \fill[ALIUSCFFFFFF] (90.951bp,-267.831bp) rectangle ++(414.960bp,-17.280bp);
  \fill[ALIUSCFFFFFF] (90.951bp,-287.511bp) rectangle ++(414.960bp,-17.040bp);
  \fill[ALIUSCFFFFFF] (90.951bp,-306.951bp) rectangle ++(414.960bp,-17.040bp);
  \fill[ALIUSCFFFFFF] (90.951bp,-326.391bp) rectangle ++(346.560bp,-17.040bp);
  \fill[ALIUSCFFFFFF] (90.951bp,-357.831bp) rectangle ++(414.960bp,-15.120bp);
  \fill[ALIUSCFFFFFF] (90.951bp,-375.111bp) rectangle ++(414.960bp,-15.120bp);
  \fill[ALIUSCFFFFFF] (90.951bp,-392.391bp) rectangle ++(414.960bp,-15.120bp);
  \fill[ALIUSCFFFFFF] (90.951bp,-409.671bp) rectangle ++(414.960bp,-15.120bp);
  \fill[ALIUSCFFFFFF] (90.951bp,-426.951bp) rectangle ++(414.960bp,-15.120bp);
  \fill[ALIUSCFFFFFF] (90.951bp,-444.231bp) rectangle ++(414.960bp,-15.120bp);
  \fill[ALIUSCFFFFFF] (90.951bp,-461.511bp) rectangle ++(414.960bp,-15.120bp);
  \fill[ALIUSCFFFFFF] (90.951bp,-478.551bp) rectangle ++(414.960bp,-15.120bp);
  \fill[ALIUSCFFFFFF] (90.951bp,-495.831bp) rectangle ++(414.960bp,-15.120bp);
  \fill[ALIUSCFFFFFF] (90.951bp,-513.111bp) rectangle ++(414.960bp,-15.120bp);
  \fill[ALIUSCFFFFFF] (90.951bp,-530.391bp) rectangle ++(414.960bp,-15.120bp);
  \fill[ALIUSCFFFFFF] (90.951bp,-547.671bp) rectangle ++(414.960bp,-15.120bp);
  \fill[ALIUSCFFFFFF] (90.951bp,-564.951bp) rectangle ++(269.280bp,-15.120bp);
  \fill[ALIUSCFFFFFF] (90.951bp,-594.231bp) rectangle ++(414.960bp,-17.040bp);
  \fill[ALIUSCFFFFFF] (90.951bp,-613.671bp) rectangle ++(414.960bp,-17.040bp);
  \fill[ALIUSCFFFFFF] (90.951bp,-633.111bp) rectangle ++(414.960bp,-17.040bp);
  \fill[ALIUSCFFFFFF] (90.951bp,-652.551bp) rectangle ++(414.960bp,-17.280bp);
  \fill[ALIUSCFFFFFF] (90.951bp,-672.231bp) rectangle ++(414.960bp,-17.040bp);
  \fill[ALIUSCFFFFFF] (90.951bp,-691.671bp) rectangle ++(414.960bp,-17.040bp);
  \fill[ALIUSCFFFFFF] (90.951bp,-711.111bp) rectangle ++(414.960bp,-17.040bp);
  \fill[ALIUSCFFFFFF] (90.951bp,-730.551bp) rectangle ++(414.960bp,-17.280bp);
  \fill[ALIUSCFFFFFF] (90.951bp,-750.231bp) rectangle ++(414.960bp,-17.040bp);
  \ALIUSPlacedTextContent{90.951}{36.329}{322.226}{ALIUSC767171}{\ALIUSFontCormorantRegular}{13.920}{Benjamin Baird – Internally generated conscious activity}
  \ALIUSPlacedTextContent{494.631}{36.329}{14.611}{ALIUSC000000}{\ALIUSFontCormorantRegular}{13.920}{25}
  \ALIUSPlacedTextContent{90.951}{73.048}{418.224}{ALIUSC000000}{\ALIUSFontCormorantRegular}{13.920}{natural systems, how the capacity for explicit self-awareness arises in the}
  \ALIUSPlacedTextContent{90.951}{92.488}{418.238}{ALIUSC000000}{\ALIUSFontCormorantRegular}{13.920}{human brain, the significance of language and symbolic representation in}
  \ALIUSPlacedTextContent{90.951}{111.928}{418.258}{ALIUSC000000}{\ALIUSFontCormorantRegular}{13.920}{defining human cognition, as well as issues pertaining to emergence,}
  \ALIUSPlacedTextContent{90.951}{131.609}{418.196}{ALIUSC000000}{\ALIUSFontCormorantRegular}{13.920}{information, and semiotic processes in biological systems. My research to}
  \ALIUSPlacedTextContent{90.951}{151.049}{418.235}{ALIUSC000000}{\ALIUSFontCormorantRegular}{13.920}{date has mostly focused on spontaneous or self-generated conscious}
  \ALIUSPlacedTextContent{90.951}{170.488}{418.164}{ALIUSC000000}{\ALIUSFontCormorantRegular}{13.920}{phenomena (e.g., mind-wandering, dreaming). I believe that studies of these}
  \ALIUSPlacedTextContent{90.951}{189.928}{418.271}{ALIUSC000000}{\ALIUSFontCormorantRegular}{13.920}{spontaneous states can provide unique and useful ways of approaching}
  \ALIUSPlacedTextContent{90.951}{209.609}{418.374}{ALIUSC000000}{\ALIUSFontCormorantRegular}{13.920}{scientific questions regarding consciousness, including characterizing its}
  \ALIUSPlacedTextContent{90.951}{229.049}{418.275}{ALIUSC000000}{\ALIUSFontCormorantRegular}{13.920}{dynamics as well as specifying the necessary conditions for its presence or}
  \ALIUSPlacedTextContent{90.951}{248.489}{418.281}{ALIUSC000000}{\ALIUSFontCormorantRegular}{13.920}{absence. Broadly, I am particularly interested in self-reflective awareness,}
  \ALIUSPlacedTextContent{90.951}{267.928}{414.974}{ALIUSC000000}{\ALIUSFontCormorantRegular}{13.920}{both the fundamental question of how a creature or entity can have self-}
  \ALIUSPlacedTextContent{90.951}{287.608}{418.259}{ALIUSC000000}{\ALIUSFontCormorantRegular}{13.920}{reflective awareness at all, as well as more generally how we monitor our}
  \ALIUSPlacedTextContent{90.951}{307.048}{418.151}{ALIUSC000000}{\ALIUSFontCormorantRegular}{13.920}{conscious states. Along these lines, I have a particular interest in lucid}
  \ALIUSPlacedTextContent{90.951}{326.488}{353.208}{ALIUSC000000}{\ALIUSFontCormorantRegular}{13.920}{dreaming, and I have focused on this state in my latest research.}
  \ALIUSPlacedTextContent{90.951}{358.053}{418.511}{ALIUSC1F8135}{\ALIUSFontLatoLight}{12.357}{Dreaming and mind wandering have been proposed to rely on similar}
  \ALIUSPlacedTextContent{90.951}{375.333}{418.389}{ALIUSC1F8135}{\ALIUSFontLatoLight}{12.357}{mechanisms to the extent that dreaming has been cast as “an intensified form}
  \ALIUSPlacedTextContent{90.951}{392.613}{418.399}{ALIUSC1F8135}{\ALIUSFontLatoLight}{12.357}{of mind-wandering” (\ALIUSCitationLink{aliusref-baird-koroma-fox-nijeboer-2013}{Fox et al., 2013}). In this account, dreaming and mind}
  \ALIUSPlacedTextContent{90.951}{409.893}{418.375}{ALIUSC1F8135}{\ALIUSFontLatoLight}{12.357}{wandering phenomena differ in quantitative rather than qualitative terms.}
  \ALIUSPlacedTextContent{90.951}{427.173}{418.292}{ALIUSC1F8135}{\ALIUSFontLatoLight}{12.357}{Supporting this proposal, electroencephalogram (EEG) measurements have}
  \ALIUSPlacedTextContent{90.951}{444.453}{418.429}{ALIUSC1F8135}{\ALIUSFontLatoLight}{12.357}{found that slow waves that are classically observed during sleep are also linked}
  \ALIUSPlacedTextContent{90.951}{461.733}{418.463}{ALIUSC1F8135}{\ALIUSFontLatoLight}{12.357}{to mind-wandering experiences during wakefulness (\ALIUSCitationLink{aliusref-baird-koroma-andrillon-windt-2019}{Andrillon et al., 2019}).}
  \ALIUSPlacedTextContent{90.951}{478.773}{418.310}{ALIUSC1F8135}{\ALIUSFontLatoLight}{12.357}{These results blur the lines between cognitive and neural processes observed}
  \ALIUSPlacedTextContent{90.951}{496.053}{418.377}{ALIUSC1F8135}{\ALIUSFontLatoLight}{12.357}{across sleep and wakefulness. Could you please explain some of the}
  \ALIUSPlacedTextContent{90.951}{513.333}{414.986}{ALIUSC1F8135}{\ALIUSFontLatoLight}{12.357}{similarities and differences between dreaming during sleep and mind-}
  \ALIUSPlacedTextContent{90.951}{530.613}{418.462}{ALIUSC1F8135}{\ALIUSFontLatoLight}{12.357}{wandering during wakefulness? Do you think we should study mind}
  \ALIUSPlacedTextContent{90.951}{547.893}{418.369}{ALIUSC1F8135}{\ALIUSFontLatoLight}{12.357}{wandering and dreaming beyond the wakefulness vs. sleep dichotomy? What}
  \ALIUSPlacedTextContent{90.951}{565.173}{272.783}{ALIUSC1F8135}{\ALIUSFontLatoLight}{12.357}{would be the advantages and pitfalls of doing so?}
  \ALIUSPlacedTextContent{90.951}{594.328}{418.423}{ALIUSC000000}{\ALIUSFontCormorantRegular}{13.920}{I think “similarities and differences” is a good way to look at it, rather than}
  \ALIUSPlacedTextContent{90.951}{613.768}{418.364}{ALIUSC000000}{\ALIUSFontCormorantRegular}{13.920}{“are they the same”? They are clearly not the same. They are both sections of}
  \ALIUSPlacedTextContent{90.951}{633.208}{418.261}{ALIUSC000000}{\ALIUSFontCormorantRegular}{13.920}{the brain-mind state-space if you like, so in that sense continuous, but there}
  \ALIUSPlacedTextContent{90.951}{652.648}{418.241}{ALIUSC000000}{\ALIUSFontCormorantRegular}{13.920}{are important differences between dreaming and mind-wandering. First, as}
  \ALIUSPlacedTextContent{90.951}{672.328}{418.271}{ALIUSC000000}{\ALIUSFontCormorantRegular}{13.920}{a number of researchers have pointed out, dreams, particularly those}
  \ALIUSPlacedTextContent{90.951}{691.768}{418.281}{ALIUSC000000}{\ALIUSFontCormorantRegular}{13.920}{occurring during Rapid Eye Movement (REM) sleep, often involve what}
  \ALIUSPlacedTextContent{90.951}{711.208}{418.275}{ALIUSC000000}{\ALIUSFontCormorantRegular}{13.920}{might be called “full immersion”, whereas mind-wandering does not. That is,}
  \ALIUSPlacedTextContent{90.951}{730.648}{418.290}{ALIUSC000000}{\ALIUSFontCormorantRegular}{13.920}{in a dream you often find yourself embodied in a dream body interacting}
  \ALIUSPlacedTextContent{90.951}{750.328}{418.243}{ALIUSC000000}{\ALIUSFontCormorantRegular}{13.920}{with a tridimensional (3D) multimodal virtual dream world. Even vivid}
  \ALIUSPlacedTextContent{90.951}{793.143}{122.397}{ALIUSC767171}{\ALIUSFontCormorantLight}{12.000}{ALIUS Bulletin n°4 (2020)}
  \ALIUSPlacedTextContent{387.435}{793.143}{120.728}{ALIUSC767171}{\ALIUSFontCormorantLight}{12.000}{aliusresearch.org/bulletin}
\end{tikzpicture}
\null
\newpage
% Page 3
\thispagestyle{empty}
\noindent\begin{tikzpicture}[remember picture,overlay,x=1bp,y=1bp,shift={(current page.north west)}]
  \fill[ALIUSCFFFFFF] (90.951bp,-72.951bp) rectangle ++(414.960bp,-17.040bp);
  \fill[ALIUSCFFFFFF] (90.951bp,-92.391bp) rectangle ++(414.960bp,-17.040bp);
  \fill[ALIUSCFFFFFF] (90.951bp,-111.831bp) rectangle ++(414.960bp,-17.280bp);
  \fill[ALIUSCFFFFFF] (90.951bp,-131.511bp) rectangle ++(414.960bp,-17.040bp);
  \fill[ALIUSCFFFFFF] (90.951bp,-150.951bp) rectangle ++(414.960bp,-17.040bp);
  \fill[ALIUSCFFFFFF] (90.951bp,-170.391bp) rectangle ++(414.960bp,-17.040bp);
  \fill[ALIUSCFFFFFF] (90.951bp,-189.831bp) rectangle ++(414.960bp,-17.280bp);
  \fill[ALIUSCFFFFFF] (90.951bp,-209.511bp) rectangle ++(414.960bp,-17.040bp);
  \fill[ALIUSCFFFFFF] (90.951bp,-228.951bp) rectangle ++(414.960bp,-17.040bp);
  \fill[ALIUSCFFFFFF] (90.951bp,-248.391bp) rectangle ++(414.960bp,-17.040bp);
  \fill[ALIUSCFFFFFF] (90.951bp,-267.831bp) rectangle ++(414.960bp,-17.280bp);
  \fill[ALIUSCFFFFFF] (90.951bp,-287.511bp) rectangle ++(216.960bp,-17.040bp);
  \fill[ALIUSCFFFFFF] (90.951bp,-326.391bp) rectangle ++(414.960bp,-17.040bp);
  \fill[ALIUSCFFFFFF] (90.951bp,-345.831bp) rectangle ++(414.960bp,-17.280bp);
  \fill[ALIUSCFFFFFF] (90.951bp,-365.511bp) rectangle ++(414.960bp,-17.040bp);
  \fill[ALIUSCFFFFFF] (90.951bp,-384.951bp) rectangle ++(414.960bp,-17.040bp);
  \fill[ALIUSCFFFFFF] (90.951bp,-404.391bp) rectangle ++(414.960bp,-17.040bp);
  \fill[ALIUSCFFFFFF] (90.951bp,-423.831bp) rectangle ++(414.960bp,-17.280bp);
  \fill[ALIUSCFFFFFF] (90.951bp,-443.511bp) rectangle ++(414.960bp,-17.040bp);
  \fill[ALIUSCFFFFFF] (90.951bp,-462.951bp) rectangle ++(414.960bp,-17.040bp);
  \fill[ALIUSCFFFFFF] (90.951bp,-482.391bp) rectangle ++(414.960bp,-17.040bp);
  \fill[ALIUSCFFFFFF] (90.951bp,-501.831bp) rectangle ++(414.960bp,-17.040bp);
  \fill[ALIUSCFFFFFF] (90.951bp,-521.271bp) rectangle ++(414.960bp,-17.280bp);
  \fill[ALIUSCFFFFFF] (90.951bp,-540.951bp) rectangle ++(414.960bp,-17.040bp);
  \fill[ALIUSCFFFFFF] (90.951bp,-560.391bp) rectangle ++(414.960bp,-17.040bp);
  \fill[ALIUSCFFFFFF] (90.951bp,-579.831bp) rectangle ++(414.960bp,-17.040bp);
  \fill[ALIUSCFFFFFF] (90.951bp,-599.271bp) rectangle ++(414.960bp,-17.280bp);
  \fill[ALIUSCFFFFFF] (90.951bp,-618.951bp) rectangle ++(414.960bp,-17.040bp);
  \fill[ALIUSCFFFFFF] (90.951bp,-638.391bp) rectangle ++(414.960bp,-17.040bp);
  \fill[ALIUSCFFFFFF] (90.951bp,-657.831bp) rectangle ++(414.960bp,-17.040bp);
  \fill[ALIUSCFFFFFF] (90.951bp,-677.271bp) rectangle ++(414.960bp,-17.280bp);
  \fill[ALIUSCFFFFFF] (90.951bp,-696.951bp) rectangle ++(382.320bp,-17.040bp);
  \fill[ALIUSCFFFFFF] (480.471bp,-696.951bp) rectangle ++(25.440bp,-17.040bp);
  \fill[ALIUSCFFFFFF] (90.951bp,-716.391bp) rectangle ++(344.880bp,-17.040bp);
  \fill[ALIUSCFFFFFF] (449.511bp,-716.391bp) rectangle ++(56.400bp,-17.040bp);
  \fill[ALIUSCFFFFFF] (90.951bp,-735.831bp) rectangle ++(414.960bp,-17.040bp);
  \ALIUSPlacedTextContent{90.951}{36.329}{322.226}{ALIUSC767171}{\ALIUSFontCormorantRegular}{13.920}{Benjamin Baird – Internally generated conscious activity}
  \ALIUSPlacedTextContent{493.671}{36.329}{15.395}{ALIUSC000000}{\ALIUSFontCormorantRegular}{13.920}{26}
  \ALIUSPlacedTextContent{90.951}{73.048}{418.234}{ALIUSC000000}{\ALIUSFontCormorantRegular}{13.920}{daydreams or mind-wandering episodes do not lead to this level of}
  \ALIUSPlacedTextContent{90.951}{92.488}{418.235}{ALIUSC000000}{\ALIUSFontCormorantRegular}{13.920}{immersiveness, or, to borrow a term from the field of virtual reality,}
  \ALIUSPlacedTextContent{90.951}{111.928}{414.974}{ALIUSC000000}{\ALIUSFontCormorantRegular}{13.920}{“presence”. From this perspective, dreaming might be regarded as a full-}
  \ALIUSPlacedTextContent{90.951}{131.609}{418.263}{ALIUSC000000}{\ALIUSFontCormorantRegular}{13.920}{blown world-simulation, which in this sense is more similar to our experience}
  \ALIUSPlacedTextContent{90.951}{151.049}{418.257}{ALIUSC000000}{\ALIUSFontCormorantRegular}{13.920}{during the waking state as a whole, rather than just mind-wandering. This}
  \ALIUSPlacedTextContent{90.951}{170.488}{418.225}{ALIUSC000000}{\ALIUSFontCormorantRegular}{13.920}{point is also related to the differences in experiential vividness of the states.}
  \ALIUSPlacedTextContent{90.951}{189.928}{418.140}{ALIUSC000000}{\ALIUSFontCormorantRegular}{13.920}{For instance, in our research we have found that individuals can perform}
  \ALIUSPlacedTextContent{90.951}{209.609}{418.329}{ALIUSC000000}{\ALIUSFontCormorantRegular}{13.920}{smooth pursuit eye movement tracking in REM sleep dreaming but not}
  \ALIUSPlacedTextContent{90.951}{229.049}{418.236}{ALIUSC000000}{\ALIUSFontCormorantRegular}{13.920}{during visuomotor imagination during wakefulness (LaBerge, Baird \&}
  \ALIUSPlacedTextContent{90.951}{248.489}{418.204}{ALIUSC000000}{\ALIUSFontCormorantRegular}{13.920}{Zimbardo, 2018). These and other findings indicate that the perceptual}
  \ALIUSPlacedTextContent{90.951}{267.928}{418.219}{ALIUSC000000}{\ALIUSFontCormorantRegular}{13.920}{vividness of dreaming is likely higher than that typically experienced during}
  \ALIUSPlacedTextContent{90.951}{287.608}{223.557}{ALIUSC000000}{\ALIUSFontCormorantRegular}{13.920}{mind-wandering or daydreaming states.}
  \ALIUSPlacedTextContent{90.951}{326.488}{418.230}{ALIUSC000000}{\ALIUSFontCormorantRegular}{13.920}{In turn (there are many overlapping constructs here!), we have proposed that}
  \ALIUSPlacedTextContent{90.951}{345.928}{418.314}{ALIUSC000000}{\ALIUSFontCormorantRegular}{13.920}{this is potentially related to the amount of disconnection or decoupling from}
  \ALIUSPlacedTextContent{90.951}{365.608}{418.251}{ALIUSC000000}{\ALIUSFontCormorantRegular}{13.920}{the external environment. While we have found that mind-wandering}
  \ALIUSPlacedTextContent{90.951}{385.048}{418.250}{ALIUSC000000}{\ALIUSFontCormorantRegular}{13.920}{consistently involves sensory decoupling, as revealed by event-related}
  \ALIUSPlacedTextContent{90.951}{404.488}{418.246}{ALIUSC000000}{\ALIUSFontCormorantRegular}{13.920}{potentials (ERP), pupillometry, and cortical phase-locking for example, the}
  \ALIUSPlacedTextContent{90.951}{423.928}{418.221}{ALIUSC000000}{\ALIUSFontCormorantRegular}{13.920}{decoupling during REM sleep is more intense, which may allow for increased}
  \ALIUSPlacedTextContent{90.951}{443.608}{418.171}{ALIUSC000000}{\ALIUSFontCormorantRegular}{13.920}{vividness due to reduced competition from external stimuli. From the}
  \ALIUSPlacedTextContent{90.951}{463.048}{418.260}{ALIUSC000000}{\ALIUSFontCormorantRegular}{13.920}{neuroscience perspective, an argument has been made that the states should}
  \ALIUSPlacedTextContent{90.951}{482.488}{418.217}{ALIUSC000000}{\ALIUSFontCormorantRegular}{13.920}{be regarded as qualitatively the same because of overlapping neural}
  \ALIUSPlacedTextContent{90.951}{501.928}{418.303}{ALIUSC000000}{\ALIUSFontCormorantRegular}{13.920}{substrates. However, while the evidence indicates that there are some brain}
  \ALIUSPlacedTextContent{90.951}{521.369}{418.238}{ALIUSC000000}{\ALIUSFontCormorantRegular}{13.920}{regions that are shared between waking mind-wandering and REM sleep}
  \ALIUSPlacedTextContent{90.951}{541.049}{418.195}{ALIUSC000000}{\ALIUSFontCormorantRegular}{13.920}{dreaming, there are also many differences, and overall in my view the}
  \ALIUSPlacedTextContent{90.951}{560.489}{418.321}{ALIUSC000000}{\ALIUSFontCormorantRegular}{13.920}{substrate of REM sleep dreaming appears to be notably different, both in}
  \ALIUSPlacedTextContent{90.951}{579.928}{418.289}{ALIUSC000000}{\ALIUSFontCormorantRegular}{13.920}{terms of neural activation patterns and neurochemistry, from that of waking}
  \ALIUSPlacedTextContent{90.951}{599.369}{418.326}{ALIUSC000000}{\ALIUSFontCormorantRegular}{13.920}{mind-wandering. Finally, part of the answer to this question hinges on how}
  \ALIUSPlacedTextContent{90.951}{619.048}{418.232}{ALIUSC000000}{\ALIUSFontCormorantRegular}{13.920}{we define mind-wandering. Jennifer Windt has made the important point}
  \ALIUSPlacedTextContent{90.951}{638.489}{418.271}{ALIUSC000000}{\ALIUSFontCormorantRegular}{13.920}{that on a broad definition of mind-wandering as spontaneous conscious}
  \ALIUSPlacedTextContent{90.951}{657.928}{418.356}{ALIUSC000000}{\ALIUSFontCormorantRegular}{13.920}{thoughts, such types of thoughts can occur in a dream! If one thing can occur}
  \ALIUSPlacedTextContent{90.951}{677.369}{418.207}{ALIUSC000000}{\ALIUSFontCormorantRegular}{13.920}{within, or during, the overall state of the other, then it seems problematic to}
  \ALIUSPlacedTextContent{90.951}{697.048}{418.252}{ALIUSC000000}{\ALIUSFontCormorantRegular}{13.920}{equate the two. There are many interesting similarities as well – i.e.,}
  \ALIUSPlacedTextContent{90.951}{716.489}{418.257}{ALIUSC000000}{\ALIUSFontCormorantRegular}{13.920}{perceptual decoupling, spontaneous internal generation, etc. – and thus}
  \ALIUSPlacedTextContent{90.951}{735.929}{418.257}{ALIUSC000000}{\ALIUSFontCormorantRegular}{13.920}{overall I find the similarities and differences approach to be the most fruitful.}
  \ALIUSPlacedTextContent{90.951}{793.143}{122.397}{ALIUSC767171}{\ALIUSFontCormorantLight}{12.000}{ALIUS Bulletin n°4 (2020)}
  \ALIUSPlacedTextContent{387.435}{793.143}{120.728}{ALIUSC767171}{\ALIUSFontCormorantLight}{12.000}{aliusresearch.org/bulletin}
\end{tikzpicture}
\null
\newpage
% Page 4
\thispagestyle{empty}
\noindent\begin{tikzpicture}[remember picture,overlay,x=1bp,y=1bp,shift={(current page.north west)}]
  \ALIUSPlacedTextContent{90.951}{36.329}{322.226}{ALIUSC767171}{\ALIUSFontCormorantRegular}{13.920}{Benjamin Baird – Internally generated conscious activity}
  \ALIUSPlacedTextContent{494.391}{36.329}{14.891}{ALIUSC000000}{\ALIUSFontCormorantRegular}{13.920}{27}
  \ALIUSPlacedTextContent{122.041}{81.307}{14.160}{ALIUSC595959}{\ALIUSFontCormorantLight}{48.000}{\ALIUSPullQuoteOpen}
  \ALIUSPlacedTextContent{140.060}{81.307}{314.879}{ALIUSC000000}{\ALIUSFontLatoLight}{15.120}{Dreams, particularly those occurring during}
  \ALIUSPlacedTextContent{140.041}{99.307}{314.919}{ALIUSC000000}{\ALIUSFontLatoLight}{15.120}{Rapid Eye Movement (REM) sleep, often involve}
  \ALIUSPlacedTextContent{140.067}{117.307}{314.865}{ALIUSC000000}{\ALIUSFontLatoLight}{15.120}{what might be called “full immersion”, whereas}
  \ALIUSPlacedTextContent{458.959}{135.307}{14.160}{ALIUSC595959}{\ALIUSFontCormorantLight}{48.000}{\ALIUSPullQuoteClose}
  \ALIUSPlacedTextContent{212.790}{135.307}{169.421}{ALIUSC000000}{\ALIUSFontLatoLight}{15.120}{mind-wandering does not.}
  \ALIUSPlacedTextContent{90.951}{181.173}{418.352}{ALIUSC1F8135}{\ALIUSFontLatoLight}{12.357}{Lucid dreaming is defined as a kind of dream during which dreamers are aware}
  \ALIUSPlacedTextContent{90.951}{198.453}{418.348}{ALIUSC1F8135}{\ALIUSFontLatoLight}{12.357}{that they are dreaming. Stephen LaBerge provided the first experimental}
  \ALIUSPlacedTextContent{90.951}{215.733}{418.395}{ALIUSC1F8135}{\ALIUSFontLatoLight}{12.357}{evidence for lucid dreaming by showing that lucid dreamers can perform in}
  \ALIUSPlacedTextContent{90.951}{233.013}{418.447}{ALIUSC1F8135}{\ALIUSFontLatoLight}{12.357}{their dream an eye-movement sequence defined with the experimenter}
  \ALIUSPlacedTextContent{90.951}{250.053}{418.345}{ALIUSC1F8135}{\ALIUSFontLatoLight}{12.357}{during wakefulness (LaBerge et al., 1981). You published recently a major}
  \ALIUSPlacedTextContent{90.951}{267.333}{418.402}{ALIUSC1F8135}{\ALIUSFontLatoLight}{12.357}{synthesis on the advances of the neurocognitive research investigating the}
  \ALIUSPlacedTextContent{90.951}{284.613}{418.395}{ALIUSC1F8135}{\ALIUSFontLatoLight}{12.357}{cerebral bases of lucid dreaming (Baird, Mota-Rolim, et al., 2019). Can you}
  \ALIUSPlacedTextContent{90.951}{301.893}{418.410}{ALIUSC1F8135}{\ALIUSFontLatoLight}{12.357}{summarize the state of what we know about the neural mechanisms of lucid}
  \ALIUSPlacedTextContent{90.951}{319.173}{55.401}{ALIUSC1F8135}{\ALIUSFontLatoLight}{12.357}{dreaming?}
  \ALIUSPlacedTextContent{90.951}{355.768}{418.302}{ALIUSC000000}{\ALIUSFontCormorantRegular}{13.920}{First, I would like to emphasize that research on this topic is still in its}
  \ALIUSPlacedTextContent{90.951}{375.208}{418.226}{ALIUSC000000}{\ALIUSFontCormorantRegular}{13.920}{infancy and we need substantially more research before definitive}
  \ALIUSPlacedTextContent{90.951}{394.888}{418.213}{ALIUSC000000}{\ALIUSFontCormorantRegular}{13.920}{conclusions can be drawn. After objectively validating lucid dreaming as a}
  \ALIUSPlacedTextContent{90.951}{414.328}{418.300}{ALIUSC000000}{\ALIUSFontCormorantRegular}{13.920}{phenomenon of REM sleep using the eye signaling method, LaBerge and}
  \ALIUSPlacedTextContent{90.951}{433.768}{418.210}{ALIUSC000000}{\ALIUSFontCormorantRegular}{13.920}{colleagues went on to study physiological correlates of lucid REM sleep using}
  \ALIUSPlacedTextContent{90.951}{453.208}{418.255}{ALIUSC000000}{\ALIUSFontCormorantRegular}{13.920}{this technique. Their team did fantastic and rigorous research, which has}
  \ALIUSPlacedTextContent{90.951}{472.888}{418.224}{ALIUSC000000}{\ALIUSFontCormorantRegular}{13.920}{unfortunately often been overlooked and not received the recognition that it}
  \ALIUSPlacedTextContent{90.951}{492.328}{418.271}{ALIUSC000000}{\ALIUSFontCormorantRegular}{13.920}{deserves. One of the main findings of that initial period of research was that}
  \ALIUSPlacedTextContent{90.951}{511.768}{418.178}{ALIUSC000000}{\ALIUSFontCormorantRegular}{13.920}{lucid REM sleep was associated with measures of phasic activation and}
  \ALIUSPlacedTextContent{90.951}{531.208}{418.287}{ALIUSC000000}{\ALIUSFontCormorantRegular}{13.920}{autonomic nervous system arousal, including higher REM density,}
  \ALIUSPlacedTextContent{90.951}{550.648}{418.214}{ALIUSC000000}{\ALIUSFontCormorantRegular}{13.920}{respiration rate and heart rate (LaBerge et al., 1986). Together these findings}
  \ALIUSPlacedTextContent{90.951}{570.328}{418.243}{ALIUSC000000}{\ALIUSFontCormorantRegular}{13.920}{suggest that lucid dreams tend to occur during periods of heightened}
  \ALIUSPlacedTextContent{90.951}{589.768}{287.563}{ALIUSC000000}{\ALIUSFontCormorantRegular}{13.920}{physiological activation during ongoing REM sleep.}
  \ALIUSPlacedTextContent{90.951}{628.648}{418.188}{ALIUSC000000}{\ALIUSFontCormorantRegular}{13.920}{These findings also raised the further question of whether the lucid REM}
  \ALIUSPlacedTextContent{90.951}{648.328}{418.319}{ALIUSC000000}{\ALIUSFontCormorantRegular}{13.920}{sleep state was associated with a global and non-specific activation of the}
  \ALIUSPlacedTextContent{90.951}{667.768}{418.229}{ALIUSC000000}{\ALIUSFontCormorantRegular}{13.920}{brain or whether it was associated with activation of specific localized brain}
  \ALIUSPlacedTextContent{90.951}{687.208}{418.217}{ALIUSC000000}{\ALIUSFontCormorantRegular}{13.920}{areas or changes in specific neural oscillatory patterns. Over the past 40 years}
  \ALIUSPlacedTextContent{90.951}{706.648}{418.275}{ALIUSC000000}{\ALIUSFontCormorantRegular}{13.920}{there have unfortunately only been about four EEG studies on lucid dreaming}
  \ALIUSPlacedTextContent{90.951}{726.328}{418.261}{ALIUSC000000}{\ALIUSFontCormorantRegular}{13.920}{published in peer-reviewed journals. And even more unfortunately, each}
  \ALIUSPlacedTextContent{90.951}{745.768}{418.258}{ALIUSC000000}{\ALIUSFontCormorantRegular}{13.920}{finds a different purported neural signature of lucid dreaming! Many of these}
  \ALIUSPlacedTextContent{90.951}{793.143}{122.397}{ALIUSC767171}{\ALIUSFontCormorantLight}{12.000}{ALIUS Bulletin n°4 (2020)}
  \ALIUSPlacedTextContent{387.435}{793.143}{120.728}{ALIUSC767171}{\ALIUSFontCormorantLight}{12.000}{aliusresearch.org/bulletin}
\end{tikzpicture}
\null
\newpage
% Page 5
\thispagestyle{empty}
\noindent\begin{tikzpicture}[remember picture,overlay,x=1bp,y=1bp,shift={(current page.north west)}]
  \ALIUSPlacedTextContent{90.951}{36.329}{322.226}{ALIUSC767171}{\ALIUSFontCormorantRegular}{13.920}{Benjamin Baird – Internally generated conscious activity}
  \ALIUSPlacedTextContent{493.431}{36.329}{15.731}{ALIUSC000000}{\ALIUSFontCormorantRegular}{13.920}{28}
  \ALIUSPlacedTextContent{90.951}{73.048}{418.315}{ALIUSC000000}{\ALIUSFontCormorantRegular}{13.920}{studies have substantial interpretive issues and limited spatial sampling of}
  \ALIUSPlacedTextContent{90.951}{92.488}{418.250}{ALIUSC000000}{\ALIUSFontCormorantRegular}{13.920}{the scalp, making them hard to interpret or compare. Thus in terms of EEG}
  \ALIUSPlacedTextContent{90.951}{111.928}{418.081}{ALIUSC000000}{\ALIUSFontCormorantRegular}{13.920}{all we can say right now is that the current research consists of mixed results}
  \ALIUSPlacedTextContent{90.951}{131.609}{418.263}{ALIUSC000000}{\ALIUSFontCormorantRegular}{13.920}{and more research is needed. In terms of functional magnetic resonance}
  \ALIUSPlacedTextContent{90.951}{151.049}{418.257}{ALIUSC000000}{\ALIUSFontCormorantRegular}{13.920}{imaging (fMRI) experiments, research is even more scant, but several recent}
  \ALIUSPlacedTextContent{90.951}{170.488}{418.374}{ALIUSC000000}{\ALIUSFontCormorantRegular}{13.920}{studies have pointed to regions of the frontoparietal network as important}
  \ALIUSPlacedTextContent{90.951}{189.928}{418.257}{ALIUSC000000}{\ALIUSFontCormorantRegular}{13.920}{for lucid dreaming (\ALIUSCitationLink{aliusref-baird-koroma-dresler-wehrle-2012}{Dresler et al., 2012}; \ALIUSCitationLink{aliusref-baird-koroma-baird-castelnovo-2018}{Baird et al., 2018}). Critically, however,}
  \ALIUSPlacedTextContent{90.951}{209.609}{418.261}{ALIUSC000000}{\ALIUSFontCormorantRegular}{13.920}{this comes from a case report and an individual differences study and there}
  \ALIUSPlacedTextContent{90.951}{229.049}{146.473}{ALIUSC000000}{\ALIUSFontCormorantRegular}{13.920}{is still no group-level fMRI}
  \ALIUSPlacedTextContent{237.462}{229.049}{6.015}{ALIUSCFF0000}{\ALIUSFontCormorantRegular}{13.920}{.}
  \ALIUSPlacedTextContent{243.251}{229.049}{265.970}{ALIUSC000000}{\ALIUSFontCormorantRegular}{13.920}{~As such, that remains one of the most important}
  \ALIUSPlacedTextContent{90.951}{248.489}{418.265}{ALIUSC000000}{\ALIUSFontCormorantRegular}{13.920}{goals for upcoming work. One of the most important findings we have comes}
  \ALIUSPlacedTextContent{90.951}{267.928}{418.242}{ALIUSC000000}{\ALIUSFontCormorantRegular}{13.920}{from pharmacology. Specifically, we know that the probability of having a}
  \ALIUSPlacedTextContent{90.951}{287.608}{418.364}{ALIUSC000000}{\ALIUSFontCormorantRegular}{13.920}{lucid dream is substantially enhanced by cholinergic stimulation during}
  \ALIUSPlacedTextContent{90.951}{307.048}{418.199}{ALIUSC000000}{\ALIUSFontCormorantRegular}{13.920}{REM sleep (LaBerge, LaMarca \& \ALIUSCitationLink{aliusref-baird-koroma-baird-castelnovo-2018}{Baird, 2018}). This fits overall with the}
  \ALIUSPlacedTextContent{90.951}{326.488}{418.257}{ALIUSC000000}{\ALIUSFontCormorantRegular}{13.920}{findings noted above suggesting lucidity is associated with increased}
  \ALIUSPlacedTextContent{90.951}{345.928}{418.212}{ALIUSC000000}{\ALIUSFontCormorantRegular}{13.920}{activation. In my view one of the most interesting and important next steps}
  \ALIUSPlacedTextContent{90.951}{365.608}{418.239}{ALIUSC000000}{\ALIUSFontCormorantRegular}{13.920}{for research on the neurobiology of lucid dreaming therefore is to understand}
  \ALIUSPlacedTextContent{90.951}{385.048}{418.239}{ALIUSC000000}{\ALIUSFontCormorantRegular}{13.920}{mechanistically why Acetylcholinesterase inhibitors have such a dramatic}
  \ALIUSPlacedTextContent{90.951}{404.488}{149.610}{ALIUSC000000}{\ALIUSFontCormorantRegular}{13.920}{effect on lucid REM sleep.}
  \ALIUSPlacedTextContent{90.951}{443.733}{418.320}{ALIUSC1F8135}{\ALIUSFontLatoLight}{12.357}{Extending the approach of LaBerge, several labs around the world recently}
  \ALIUSPlacedTextContent{90.951}{460.773}{418.495}{ALIUSC1F8135}{\ALIUSFontLatoLight}{12.357}{investigated the ability of experimenters to communicate with the lucid}
  \ALIUSPlacedTextContent{90.951}{478.053}{418.312}{ALIUSC1F8135}{\ALIUSFontLatoLight}{12.357}{dreamer. To do so, they observed the response of the lucid dreamer to}
  \ALIUSPlacedTextContent{90.951}{495.333}{418.340}{ALIUSC1F8135}{\ALIUSFontLatoLight}{12.357}{questions (e.g., arithmetic operations) using a variety of predefined “codes”,}
  \ALIUSPlacedTextContent{90.951}{512.613}{418.387}{ALIUSC1F8135}{\ALIUSFontLatoLight}{12.357}{such as moving the eyes in one direction or moving certain muscles of the}
  \ALIUSPlacedTextContent{90.951}{529.893}{418.330}{ALIUSC1F8135}{\ALIUSFontLatoLight}{12.357}{face to say “yes” and moving the eyes in another direction or moving other}
  \ALIUSPlacedTextContent{90.951}{547.173}{418.401}{ALIUSC1F8135}{\ALIUSFontLatoLight}{12.357}{muscles for “no”. Which new insights can be gained by the development of an}
  \ALIUSPlacedTextContent{90.951}{564.453}{313.244}{ALIUSC1F8135}{\ALIUSFontLatoLight}{12.357}{extended real-time communication with a lucid dreamer?}
  \ALIUSPlacedTextContent{90.951}{593.609}{418.224}{ALIUSC000000}{\ALIUSFontCormorantRegular}{13.920}{I think these findings are important to publish in the peer-reviewed literature}
  \ALIUSPlacedTextContent{90.951}{613.048}{418.140}{ALIUSC000000}{\ALIUSFontCormorantRegular}{13.920}{as a proof of concept. Nevertheless, it is perhaps worth mentioning that many}
  \ALIUSPlacedTextContent{90.951}{632.489}{418.201}{ALIUSC000000}{\ALIUSFontCormorantRegular}{13.920}{of us in the field have known for a long time that this was possible since there}
  \ALIUSPlacedTextContent{90.951}{651.928}{418.207}{ALIUSC000000}{\ALIUSFontCormorantRegular}{13.920}{have been lucid dream induction devices available for decades that include}
  \ALIUSPlacedTextContent{90.951}{671.369}{418.244}{ALIUSC000000}{\ALIUSFontCormorantRegular}{13.920}{the possibility for two-way interaction with the device through eye}
  \ALIUSPlacedTextContent{90.951}{691.048}{418.264}{ALIUSC000000}{\ALIUSFontCormorantRegular}{13.920}{movements. I haven’t yet heard a compelling argument for how this will open}
  \ALIUSPlacedTextContent{90.951}{710.489}{418.118}{ALIUSC000000}{\ALIUSFontCormorantRegular}{13.920}{up new avenues of research. Most of the time lucid dreamers should be in a}
  \ALIUSPlacedTextContent{90.951}{729.929}{418.239}{ALIUSC000000}{\ALIUSFontCormorantRegular}{13.920}{position to remember experimental tasks and intentions that were set in}
  \ALIUSPlacedTextContent{90.951}{749.369}{418.265}{ALIUSC000000}{\ALIUSFontCormorantRegular}{13.920}{conversation with an experimenter before going to sleep, so it is hard to see}
  \ALIUSPlacedTextContent{90.951}{793.143}{122.397}{ALIUSC767171}{\ALIUSFontCormorantLight}{12.000}{ALIUS Bulletin n°4 (2020)}
  \ALIUSPlacedTextContent{387.435}{793.143}{120.728}{ALIUSC767171}{\ALIUSFontCormorantLight}{12.000}{aliusresearch.org/bulletin}
\end{tikzpicture}
\null
\newpage
% Page 6
\thispagestyle{empty}
\noindent\begin{tikzpicture}[remember picture,overlay,x=1bp,y=1bp,shift={(current page.north west)}]
  \ALIUSPlacedTextContent{90.951}{36.329}{322.226}{ALIUSC767171}{\ALIUSFontCormorantRegular}{13.920}{Benjamin Baird – Internally generated conscious activity}
  \ALIUSPlacedTextContent{493.671}{36.329}{15.395}{ALIUSC000000}{\ALIUSFontCormorantRegular}{13.920}{29}
  \ALIUSPlacedTextContent{90.951}{73.048}{418.272}{ALIUSC000000}{\ALIUSFontCormorantRegular}{13.920}{what this adds that goes beyond that in a substantially different way.}
  \ALIUSPlacedTextContent{90.951}{92.488}{418.325}{ALIUSC000000}{\ALIUSFontCormorantRegular}{13.920}{Furthermore, communicating in this way with someone who is asleep and}
  \ALIUSPlacedTextContent{90.951}{111.928}{418.192}{ALIUSC000000}{\ALIUSFontCormorantRegular}{13.920}{dreaming consistently runs the risk of waking him or her up. These comments}
  \ALIUSPlacedTextContent{90.951}{131.609}{418.296}{ALIUSC000000}{\ALIUSFontCormorantRegular}{13.920}{should not in any way be taken to denigrate this research, as again I think it}
  \ALIUSPlacedTextContent{90.951}{151.049}{418.303}{ALIUSC000000}{\ALIUSFontCormorantRegular}{13.920}{is important to publish as proof-of-concept. And perhaps there could be uses}
  \ALIUSPlacedTextContent{90.951}{170.488}{418.225}{ALIUSC000000}{\ALIUSFontCormorantRegular}{13.920}{for this that I have not thought of yet. On the other hand, it does strike me}
  \ALIUSPlacedTextContent{90.951}{189.928}{418.300}{ALIUSC000000}{\ALIUSFontCormorantRegular}{13.920}{that having a higher repertoire of signals on the output side could be useful}
  \ALIUSPlacedTextContent{90.951}{209.609}{418.299}{ALIUSC000000}{\ALIUSFontCormorantRegular}{13.920}{for “annotating” ongoing physiological sleep recordings. For instance, my}
  \ALIUSPlacedTextContent{90.951}{229.049}{418.190}{ALIUSC000000}{\ALIUSFontCormorantRegular}{13.920}{colleague Stephen LaBerge has suggested that by using something like 10}
  \ALIUSPlacedTextContent{90.951}{248.489}{418.245}{ALIUSC000000}{\ALIUSFontCormorantRegular}{13.920}{different types of eye signals a dreamer in the sleep lab could indicate that}
  \ALIUSPlacedTextContent{90.951}{267.928}{418.272}{ALIUSC000000}{\ALIUSFontCormorantRegular}{13.920}{they are currently experiencing a variety of specific kinds of conscious}
  \ALIUSPlacedTextContent{90.951}{287.608}{418.265}{ALIUSC000000}{\ALIUSFontCormorantRegular}{13.920}{content. By making these distinct signals, they could in real time label}
  \ALIUSPlacedTextContent{90.951}{307.048}{418.303}{ALIUSC000000}{\ALIUSFontCormorantRegular}{13.920}{ongoing physiological recordings of dreams. This could be highly useful for}
  \ALIUSPlacedTextContent{90.951}{326.488}{382.410}{ALIUSC000000}{\ALIUSFontCormorantRegular}{13.920}{psychophysiological studies of mind-brain relationships during sleep.}
  \ALIUSPlacedTextContent{90.951}{358.053}{418.281}{ALIUSC1F8135}{\ALIUSFontLatoLight}{12.357}{Lucid dreaming differs from dreaming in terms of metacognition, the cognitive}
  \ALIUSPlacedTextContent{90.951}{375.333}{418.273}{ALIUSC1F8135}{\ALIUSFontLatoLight}{12.357}{ability of having explicit knowledge on our cognitive processes. Indeed, we}
  \ALIUSPlacedTextContent{90.951}{392.613}{418.414}{ALIUSC1F8135}{\ALIUSFontLatoLight}{12.357}{know that we are dreaming when we are dreaming lucidly as contrary to}
  \ALIUSPlacedTextContent{90.951}{409.893}{418.466}{ALIUSC1F8135}{\ALIUSFontLatoLight}{12.357}{regular dreaming (Kahan \& LaBerge, 1994). Yet, lucid dreaming is also}
  \ALIUSPlacedTextContent{90.951}{427.173}{418.473}{ALIUSC1F8135}{\ALIUSFontLatoLight}{12.357}{characterized by other mental capacities that differ from regular dreaming.}
  \ALIUSPlacedTextContent{90.951}{444.453}{418.287}{ALIUSC1F8135}{\ALIUSFontLatoLight}{12.357}{For example, voluntary control is rated by the lucid dreamer as similar as}
  \ALIUSPlacedTextContent{90.951}{461.733}{418.399}{ALIUSC1F8135}{\ALIUSFontLatoLight}{12.357}{wakefulness and is enhanced as compared to normal dreaming (Dresler et al.,}
  \ALIUSPlacedTextContent{90.951}{478.773}{418.305}{ALIUSC1F8135}{\ALIUSFontLatoLight}{12.357}{2014). This allows for the dreamer to take the control of the course of the}
  \ALIUSPlacedTextContent{90.951}{496.053}{418.387}{ALIUSC1F8135}{\ALIUSFontLatoLight}{12.357}{dreaming scenario. What can lucid dreaming teach us about the role of}
  \ALIUSPlacedTextContent{90.951}{513.333}{418.454}{ALIUSC1F8135}{\ALIUSFontLatoLight}{12.357}{metacognition and more generally about the difference in cognitive activity}
  \ALIUSPlacedTextContent{90.951}{530.613}{200.965}{ALIUSC1F8135}{\ALIUSFontLatoLight}{12.357}{during dreams and the waking state?}
  \ALIUSPlacedTextContent{90.951}{559.768}{418.039}{ALIUSC000000}{\ALIUSFontCormorantRegular}{13.920}{You ask an interesting and important question. In fact, the paper you cite by}
  \ALIUSPlacedTextContent{90.951}{579.208}{418.234}{ALIUSC000000}{\ALIUSFontCormorantRegular}{13.920}{Kahan and LaBerge showed that non-lucid dreams are characterized by more}
  \ALIUSPlacedTextContent{90.951}{598.648}{418.060}{ALIUSC000000}{\ALIUSFontCormorantRegular}{13.920}{metacognitive thought than people typically acknowledge. For instance, they}
  \ALIUSPlacedTextContent{90.951}{618.328}{418.260}{ALIUSC000000}{\ALIUSFontCormorantRegular}{13.920}{found that people often experience high-level cognition in dreams, including}
  \ALIUSPlacedTextContent{90.951}{637.768}{418.259}{ALIUSC000000}{\ALIUSFontCormorantRegular}{13.920}{thoughts such as “I wonder how what I just said may have caused so-and-so}
  \ALIUSPlacedTextContent{90.951}{657.208}{418.340}{ALIUSC000000}{\ALIUSFontCormorantRegular}{13.920}{to feel” or “so-and-so asked me a question and I realized I didn’t know the}
  \ALIUSPlacedTextContent{90.951}{676.648}{418.252}{ALIUSC000000}{\ALIUSFontCormorantRegular}{13.920}{answer and I felt embarrassed”, etc. These are remarkably complex thoughts}
  \ALIUSPlacedTextContent{90.951}{696.328}{418.278}{ALIUSC000000}{\ALIUSFontCormorantRegular}{13.920}{that have metacognitive components i.e., directly reflecting on one’s own}
  \ALIUSPlacedTextContent{90.951}{715.768}{418.232}{ALIUSC000000}{\ALIUSFontCormorantRegular}{13.920}{cognition or emotional state and/or others (theory-of-mind). Nevertheless,}
  \ALIUSPlacedTextContent{90.951}{735.208}{418.257}{ALIUSC000000}{\ALIUSFontCormorantRegular}{13.920}{despite this, by definition we lack a specific kind of metacognitive knowledge}
  \ALIUSPlacedTextContent{90.951}{793.143}{122.397}{ALIUSC767171}{\ALIUSFontCormorantLight}{12.000}{ALIUS Bulletin n°4 (2020)}
  \ALIUSPlacedTextContent{387.435}{793.143}{120.728}{ALIUSC767171}{\ALIUSFontCormorantLight}{12.000}{aliusresearch.org/bulletin}
\end{tikzpicture}
\null
\newpage
% Page 7
\thispagestyle{empty}
\noindent\begin{tikzpicture}[remember picture,overlay,x=1bp,y=1bp,shift={(current page.north west)}]
  \ALIUSPlacedTextContent{90.951}{36.329}{322.226}{ALIUSC767171}{\ALIUSFontCormorantRegular}{13.920}{Benjamin Baird – Internally generated conscious activity}
  \ALIUSPlacedTextContent{493.911}{36.329}{15.381}{ALIUSC000000}{\ALIUSFontCormorantRegular}{13.920}{30}
  \ALIUSPlacedTextContent{90.951}{73.048}{418.173}{ALIUSC000000}{\ALIUSFontCormorantRegular}{13.920}{in (non-lucid) dreams: namely, knowledge about the overall state of}
  \ALIUSPlacedTextContent{90.951}{92.488}{418.318}{ALIUSC000000}{\ALIUSFontCormorantRegular}{13.920}{consciousness we are in. On the basis of these findings, Kahan and LaBerge}
  \ALIUSPlacedTextContent{90.951}{111.928}{418.288}{ALIUSC000000}{\ALIUSFontCormorantRegular}{13.920}{argued that dreaming and wakefulness actually aren’t that different in terms}
  \ALIUSPlacedTextContent{90.951}{131.609}{418.265}{ALIUSC000000}{\ALIUSFontCormorantRegular}{13.920}{of ongoing thought. After all, most of us don’t go around during the waking}
  \ALIUSPlacedTextContent{90.951}{151.049}{418.288}{ALIUSC000000}{\ALIUSFontCormorantRegular}{13.920}{state reflecting on our state of being awake, and so it was argued that it’s}
  \ALIUSPlacedTextContent{90.951}{170.488}{421.533}{ALIUSC000000}{\ALIUSFontCormorantRegular}{13.920}{actually not that peculiar that we do not do this during dreaming either.}
  \ALIUSPlacedTextContent{90.951}{189.928}{418.269}{ALIUSC000000}{\ALIUSFontCormorantRegular}{13.920}{However, a problem with this argument is that dreaming scenarios and}
  \ALIUSPlacedTextContent{90.951}{209.609}{418.257}{ALIUSC000000}{\ALIUSFontCormorantRegular}{13.920}{content are not merely a recapitulation of our waking lives. Instead, not only}
  \ALIUSPlacedTextContent{90.951}{229.049}{418.204}{ALIUSC000000}{\ALIUSFontCormorantRegular}{13.920}{do highly bizarre events occur during dreams, but I can experience myself as}
  \ALIUSPlacedTextContent{90.951}{248.489}{418.129}{ALIUSC000000}{\ALIUSFontCormorantRegular}{13.920}{a totally different person, or living in a totally different century, or I could}
  \ALIUSPlacedTextContent{90.951}{267.928}{418.210}{ALIUSC000000}{\ALIUSFontCormorantRegular}{13.920}{find myself on another planet! The point is that during dreams we can find}
  \ALIUSPlacedTextContent{90.951}{287.608}{418.235}{ALIUSC000000}{\ALIUSFontCormorantRegular}{13.920}{ourselves in situations that are profoundly discontinuous with the rest of our}
  \ALIUSPlacedTextContent{90.951}{307.048}{418.332}{ALIUSC000000}{\ALIUSFontCormorantRegular}{13.920}{waking lives and yet we still don’t notice that anything is amiss. This shows}
  \ALIUSPlacedTextContent{90.951}{326.488}{418.183}{ALIUSC000000}{\ALIUSFontCormorantRegular}{13.920}{how the picture is actually more complex than is typically discussed: on the}
  \ALIUSPlacedTextContent{90.951}{345.928}{418.316}{ALIUSC000000}{\ALIUSFontCormorantRegular}{13.920}{one hand, we can have highly complex, including metacognitive, thoughts,}
  \ALIUSPlacedTextContent{90.951}{365.608}{418.329}{ALIUSC000000}{\ALIUSFontCormorantRegular}{13.920}{and on the other hand, we can be completely disoriented, which suggests that}
  \ALIUSPlacedTextContent{90.951}{385.048}{418.264}{ALIUSC000000}{\ALIUSFontCormorantRegular}{13.920}{something is profoundly different about our state of consciousness during}
  \ALIUSPlacedTextContent{90.951}{404.488}{48.547}{ALIUSC000000}{\ALIUSFontCormorantRegular}{13.920}{dreams.}
  \ALIUSPlacedTextContent{90.951}{435.928}{418.328}{ALIUSC000000}{\ALIUSFontCormorantRegular}{13.920}{I think an interesting direction to pursue in resolving this apparent paradox}
  \ALIUSPlacedTextContent{90.951}{455.608}{418.157}{ALIUSC000000}{\ALIUSFontCormorantRegular}{13.920}{is Tulving’s notion of “autonoetic consciousness” (\ALIUSCitationLink{aliusref-baird-koroma-tulving-1985}{Tulving, 1985}). In other}
  \ALIUSPlacedTextContent{90.951}{475.048}{418.236}{ALIUSC000000}{\ALIUSFontCormorantRegular}{13.920}{research we find that, contrary to thoughts in the waking state and}
  \ALIUSPlacedTextContent{90.951}{494.488}{418.221}{ALIUSC000000}{\ALIUSFontCormorantRegular}{13.920}{particularly mind-wandering, dreaming individuals rarely engage in episodic}
  \ALIUSPlacedTextContent{90.951}{513.928}{418.201}{ALIUSC000000}{\ALIUSFontCormorantRegular}{13.920}{memory or autobiographical planning for the future. Thus one way to think}
  \ALIUSPlacedTextContent{90.951}{533.369}{414.974}{ALIUSC000000}{\ALIUSFontCormorantRegular}{13.920}{about what is different during dreaming consciousness is in terms of our self-}
  \ALIUSPlacedTextContent{90.951}{553.049}{418.265}{ALIUSC000000}{\ALIUSFontCormorantRegular}{13.920}{model, including our experience of ourselves as a self extended over time.}
  \ALIUSPlacedTextContent{90.951}{572.489}{418.159}{ALIUSC000000}{\ALIUSFontCormorantRegular}{13.920}{Jennifer Windt and Thomas Metzinger have done excellent work in this area,}
  \ALIUSPlacedTextContent{90.951}{591.928}{418.249}{ALIUSC000000}{\ALIUSFontCormorantRegular}{13.920}{and have begun to think about how many of the key changes in consciousness}
  \ALIUSPlacedTextContent{90.951}{611.369}{418.249}{ALIUSC000000}{\ALIUSFontCormorantRegular}{13.920}{that occur after becoming lucid that you note above – including volitional}
  \ALIUSPlacedTextContent{90.951}{631.048}{418.259}{ALIUSC000000}{\ALIUSFontCormorantRegular}{13.920}{control, episodic memory and metacognitive awareness of state – may all be}
  \ALIUSPlacedTextContent{90.951}{650.489}{355.532}{ALIUSC000000}{\ALIUSFontCormorantRegular}{13.920}{related to changes in our self-model (Windt \& Metzinger, 2007).}
  \ALIUSPlacedTextContent{119.262}{688.747}{14.160}{ALIUSC595959}{\ALIUSFontCormorantLight}{48.000}{\ALIUSPullQuoteOpen}
  \ALIUSPlacedTextContent{137.262}{688.747}{320.476}{ALIUSC000000}{\ALIUSFontLatoLight}{15.120}{One way to think about what is different during}
  \ALIUSPlacedTextContent{139.183}{706.747}{316.634}{ALIUSC000000}{\ALIUSFontLatoLight}{15.120}{dreaming consciousness is in terms of our self-}
  \ALIUSPlacedTextContent{137.300}{724.747}{320.400}{ALIUSC000000}{\ALIUSFontLatoLight}{15.120}{model, including our experience of ourselves as a}
  \ALIUSPlacedTextContent{461.738}{742.747}{14.160}{ALIUSC595959}{\ALIUSFontCormorantLight}{48.000}{\ALIUSPullQuoteClose}
  \ALIUSPlacedTextContent{219.293}{742.747}{156.413}{ALIUSC000000}{\ALIUSFontLatoLight}{15.120}{self extended over time.}
  \ALIUSPlacedTextContent{90.951}{793.143}{122.397}{ALIUSC767171}{\ALIUSFontCormorantLight}{12.000}{ALIUS Bulletin n°4 (2020)}
  \ALIUSPlacedTextContent{387.435}{793.143}{120.728}{ALIUSC767171}{\ALIUSFontCormorantLight}{12.000}{aliusresearch.org/bulletin}
\end{tikzpicture}
\null
\newpage
% Page 8
\thispagestyle{empty}
\noindent\begin{tikzpicture}[remember picture,overlay,x=1bp,y=1bp,shift={(current page.north west)}]
  \ALIUSPlacedTextContent{90.951}{36.329}{322.226}{ALIUSC767171}{\ALIUSFontCormorantRegular}{13.920}{Benjamin Baird – Internally generated conscious activity}
  \ALIUSPlacedTextContent{495.831}{36.329}{13.351}{ALIUSC000000}{\ALIUSFontCormorantRegular}{13.920}{31}
  \ALIUSPlacedTextContent{90.951}{73.048}{418.315}{ALIUSC000000}{\ALIUSFontCormorantRegular}{13.920}{Specifically, they have made the interesting suggestion that the specific}
  \ALIUSPlacedTextContent{90.951}{92.488}{418.243}{ALIUSC000000}{\ALIUSFontCormorantRegular}{13.920}{metacognitive deficit in dreams may be directly linked to our ability to think}
  \ALIUSPlacedTextContent{90.951}{111.928}{414.974}{ALIUSC000000}{\ALIUSFontCormorantRegular}{13.920}{about our self being in a current state. Similarly, as noted above, the re-}
  \ALIUSPlacedTextContent{90.951}{131.609}{418.201}{ALIUSC000000}{\ALIUSFontCormorantRegular}{13.920}{instantiation of episodic memory also seems directly linked to our ability to}
  \ALIUSPlacedTextContent{90.951}{151.049}{418.293}{ALIUSC000000}{\ALIUSFontCormorantRegular}{13.920}{experience selfhood, specifically the experience of ourselves over time, and}
  \ALIUSPlacedTextContent{90.951}{170.488}{418.247}{ALIUSC000000}{\ALIUSFontCormorantRegular}{13.920}{volitional control is the experience of myself as an agent that can direct my}
  \ALIUSPlacedTextContent{90.951}{189.928}{418.107}{ALIUSC000000}{\ALIUSFontCormorantRegular}{13.920}{attention or actions. Overall this suggests that for humans a key difference}
  \ALIUSPlacedTextContent{90.951}{209.609}{418.285}{ALIUSC000000}{\ALIUSFontCormorantRegular}{13.920}{between waking and dreaming may be in how we experience ourselves as}
  \ALIUSPlacedTextContent{90.951}{229.049}{108.213}{ALIUSC000000}{\ALIUSFontCormorantRegular}{13.920}{conscious subjects.}
  \ALIUSPlacedTextContent{90.951}{260.613}{418.443}{ALIUSC1F8135}{\ALIUSFontLatoLight}{12.357}{For inexperienced lucid dreamers, dream reports reveal that it is difficult to}
  \ALIUSPlacedTextContent{90.951}{277.893}{418.412}{ALIUSC1F8135}{\ALIUSFontLatoLight}{12.357}{maintain the state of lucidity and attempts to control the dream scenario may}
  \ALIUSPlacedTextContent{90.951}{295.173}{418.457}{ALIUSC1F8135}{\ALIUSFontLatoLight}{12.357}{lead to the loss of lucidity or to awakenings (\ALIUSCitationLink{aliusref-baird-koroma-laberge-rheingold-1990}{LaBerge \& Rheingold, 1990}). Yet,}
  \ALIUSPlacedTextContent{90.951}{312.453}{418.493}{ALIUSC1F8135}{\ALIUSFontLatoLight}{12.357}{with training, expert lucid dreamers can learn how to maintain such a state}
  \ALIUSPlacedTextContent{90.951}{329.733}{418.272}{ALIUSC1F8135}{\ALIUSFontLatoLight}{12.357}{(Stumbrys et al., 2012). Such consideration can be extended to meditation as}
  \ALIUSPlacedTextContent{90.951}{347.013}{418.252}{ALIUSC1F8135}{\ALIUSFontLatoLight}{12.357}{beginners might have difficulty entering the meditative state or remaining in}
  \ALIUSPlacedTextContent{90.951}{364.053}{418.403}{ALIUSC1F8135}{\ALIUSFontLatoLight}{12.357}{it (Brown \& Engler, 1980). How does training allow for the access and}
  \ALIUSPlacedTextContent{90.951}{381.333}{418.387}{ALIUSC1F8135}{\ALIUSFontLatoLight}{12.357}{stabilization of these conscious states? How do they relate to interindividual}
  \ALIUSPlacedTextContent{90.951}{398.613}{405.226}{ALIUSC1F8135}{\ALIUSFontLatoLight}{12.357}{traits and practices (\ALIUSCitationLink{aliusref-baird-koroma-schredl-erlacher-2004}{Schredl \& Erlacher, 2004}; Baird, Riedner, et al., 2019)?}
  \ALIUSPlacedTextContent{90.951}{427.768}{418.087}{ALIUSC000000}{\ALIUSFontCormorantRegular}{13.920}{Your question points to the connections between lucid dreaming and}
  \ALIUSPlacedTextContent{90.951}{447.208}{418.169}{ALIUSC000000}{\ALIUSFontCormorantRegular}{13.920}{meditation, which to my mind is an interesting topic deserving of more}
  \ALIUSPlacedTextContent{90.951}{466.888}{418.164}{ALIUSC000000}{\ALIUSFontCormorantRegular}{13.920}{research. There are two distinct aspects: access and stabilization. We have}
  \ALIUSPlacedTextContent{90.951}{486.328}{418.304}{ALIUSC000000}{\ALIUSFontCormorantRegular}{13.920}{found that long-term meditation practitioners report more frequent lucid}
  \ALIUSPlacedTextContent{90.951}{505.768}{418.223}{ALIUSC000000}{\ALIUSFontCormorantRegular}{13.920}{dreams compared to individuals without meditation training, suggesting that}
  \ALIUSPlacedTextContent{90.951}{525.208}{418.343}{ALIUSC000000}{\ALIUSFontCormorantRegular}{13.920}{there is something about meditation practice that leads to greater access to}
  \ALIUSPlacedTextContent{90.951}{544.648}{418.253}{ALIUSC000000}{\ALIUSFontCormorantRegular}{13.920}{the lucid dreaming state (Baird et al., 2019). Additionally, many of the}
  \ALIUSPlacedTextContent{90.951}{564.328}{414.974}{ALIUSC000000}{\ALIUSFontCormorantRegular}{13.920}{primary skills cultivated in meditation practice (particularly open-}
  \ALIUSPlacedTextContent{90.951}{583.768}{418.163}{ALIUSC000000}{\ALIUSFontCormorantRegular}{13.920}{monitoring or focused-attention meditation), including stability of attention}
  \ALIUSPlacedTextContent{90.951}{603.208}{418.242}{ALIUSC000000}{\ALIUSFontCormorantRegular}{13.920}{and meta-awareness, are thought to be useful in having lucid dreams.}
  \ALIUSPlacedTextContent{90.951}{622.648}{418.293}{ALIUSC000000}{\ALIUSFontCormorantRegular}{13.920}{Stumbrys et al. (2015) found an intriguing association between dispositional}
  \ALIUSPlacedTextContent{90.951}{642.328}{418.277}{ALIUSC000000}{\ALIUSFontCormorantRegular}{13.920}{mindfulness (a construct that broadly refers to cultivating awareness of}
  \ALIUSPlacedTextContent{90.951}{661.768}{418.393}{ALIUSC000000}{\ALIUSFontCormorantRegular}{13.920}{experience in the present moment) during waking states and lucid dream}
  \ALIUSPlacedTextContent{90.951}{681.208}{418.235}{ALIUSC000000}{\ALIUSFontCormorantRegular}{13.920}{frequency, but only in individuals with prior meditation experience.}
  \ALIUSPlacedTextContent{90.951}{700.648}{418.287}{ALIUSC000000}{\ALIUSFontCormorantRegular}{13.920}{Although preliminary, one interpretation of this finding is that at least some}
  \ALIUSPlacedTextContent{90.951}{720.328}{418.319}{ALIUSC000000}{\ALIUSFontCormorantRegular}{13.920}{types of meditation practices result in changes in trait mindfulness, or}
  \ALIUSPlacedTextContent{90.951}{739.768}{418.185}{ALIUSC000000}{\ALIUSFontCormorantRegular}{13.920}{cognitive skills associated with specific aspects of mindfulness, that then}
  \ALIUSPlacedTextContent{90.951}{793.143}{122.397}{ALIUSC767171}{\ALIUSFontCormorantLight}{12.000}{ALIUS Bulletin n°4 (2020)}
  \ALIUSPlacedTextContent{387.435}{793.143}{120.728}{ALIUSC767171}{\ALIUSFontCormorantLight}{12.000}{aliusresearch.org/bulletin}
\end{tikzpicture}
\null
\newpage
% Page 9
\thispagestyle{empty}
\noindent\begin{tikzpicture}[remember picture,overlay,x=1bp,y=1bp,shift={(current page.north west)}]
  \ALIUSPlacedTextContent{90.951}{36.329}{322.226}{ALIUSC767171}{\ALIUSFontCormorantRegular}{13.920}{Benjamin Baird – Internally generated conscious activity}
  \ALIUSPlacedTextContent{494.871}{36.329}{14.331}{ALIUSC000000}{\ALIUSFontCormorantRegular}{13.920}{32}
  \ALIUSPlacedTextContent{90.951}{73.048}{418.256}{ALIUSC000000}{\ALIUSFontCormorantRegular}{13.920}{carry over into sleep and dream states, leading to increases in lucidity. Our}
  \ALIUSPlacedTextContent{90.951}{92.488}{37.495}{ALIUSC000000}{\ALIUSFontCormorantRegular}{13.920}{recent}
  \ALIUSPlacedTextContent{152.612}{92.488}{48.202}{ALIUSC000000}{\ALIUSFontCormorantRegular}{13.920}{research}
  \ALIUSPlacedTextContent{224.981}{92.488}{20.842}{ALIUSC000000}{\ALIUSFontCormorantRegular}{13.920}{has}
  \ALIUSPlacedTextContent{269.990}{92.488}{24.171}{ALIUSC000000}{\ALIUSFontCormorantRegular}{13.920}{also}
  \ALIUSPlacedTextContent{318.327}{92.488}{35.559}{ALIUSC000000}{\ALIUSFontCormorantRegular}{13.920}{found}
  \ALIUSPlacedTextContent{378.052}{92.488}{9.119}{ALIUSC000000}{\ALIUSFontCormorantRegular}{13.920}{a}
  \ALIUSPlacedTextContent{411.337}{92.488}{24.797}{ALIUSC000000}{\ALIUSFontCormorantRegular}{13.920}{link}
  \ALIUSPlacedTextContent{460.301}{92.488}{52.183}{ALIUSC000000}{\ALIUSFontCormorantRegular}{13.920}{between}
  \ALIUSPlacedTextContent{90.951}{111.928}{421.533}{ALIUSC000000}{\ALIUSFontCormorantRegular}{13.920}{specific aspects of trait mindfulness and lucid dream frequency.}
  \ALIUSPlacedTextContent{90.951}{131.609}{418.241}{ALIUSC000000}{\ALIUSFontCormorantRegular}{13.920}{The link between mental training and sustaining lucidity is equally}
  \ALIUSPlacedTextContent{90.951}{151.049}{418.279}{ALIUSC000000}{\ALIUSFontCormorantRegular}{13.920}{interesting from a theoretical point of view, however almost no empirical}
  \ALIUSPlacedTextContent{90.951}{170.488}{418.252}{ALIUSC000000}{\ALIUSFontCormorantRegular}{13.920}{work has been done to explore this connection. Dunne, Thompson and}
  \ALIUSPlacedTextContent{90.951}{189.928}{418.261}{ALIUSC000000}{\ALIUSFontCormorantRegular}{13.920}{Schooler (2019) have recently argued that certain styles of mediation, for}
  \ALIUSPlacedTextContent{90.951}{209.609}{414.974}{ALIUSC000000}{\ALIUSFontCormorantRegular}{13.920}{instance Tibetan Mahamudra, specifically cultivate a type of non-}
  \ALIUSPlacedTextContent{90.951}{229.049}{418.234}{ALIUSC000000}{\ALIUSFontCormorantRegular}{13.920}{propositional, sustained meta-awareness that could be useful in sustaining}
  \ALIUSPlacedTextContent{90.951}{248.489}{418.268}{ALIUSC000000}{\ALIUSFontCormorantRegular}{13.920}{lucidity during a dream. I would like to see this investigated in future}
  \ALIUSPlacedTextContent{90.951}{267.928}{50.960}{ALIUSC000000}{\ALIUSFontCormorantRegular}{13.920}{research.}
  \ALIUSPlacedTextContent{90.951}{299.733}{418.403}{ALIUSC1F8135}{\ALIUSFontLatoLight}{12.357}{Sleep represents an interesting case for consciousness research since it}
  \ALIUSPlacedTextContent{90.951}{317.013}{418.519}{ALIUSC1F8135}{\ALIUSFontLatoLight}{12.357}{consists in periods during which conscious activity is reduced, typically in}
  \ALIUSPlacedTextContent{90.951}{334.053}{418.420}{ALIUSC1F8135}{\ALIUSFontLatoLight}{12.357}{NREM sleep, and a period where conscious activity qualitatively differs from}
  \ALIUSPlacedTextContent{90.951}{351.333}{418.428}{ALIUSC1F8135}{\ALIUSFontLatoLight}{12.357}{wakefulness, typically in REM sleep. Yet, conscious experiences during sleep}
  \ALIUSPlacedTextContent{90.951}{368.613}{418.320}{ALIUSC1F8135}{\ALIUSFontLatoLight}{12.357}{have long been considered with caution as reports about dreaming activity}
  \ALIUSPlacedTextContent{90.951}{385.893}{418.329}{ALIUSC1F8135}{\ALIUSFontLatoLight}{12.357}{are available only after the sleeper has awakened and not during sleep}
  \ALIUSPlacedTextContent{90.951}{403.173}{418.306}{ALIUSC1F8135}{\ALIUSFontLatoLight}{12.357}{(\ALIUSCitationLink{aliusref-baird-koroma-malcolm-1958}{Malcolm, 1958}; \ALIUSCitationLink{aliusref-baird-koroma-windt-2013}{Windt, 2013}). The identification of neural markers of}
  \ALIUSPlacedTextContent{90.951}{420.453}{418.334}{ALIUSC1F8135}{\ALIUSFontLatoLight}{12.357}{consciousness, as well as the development of lucid dreaming research, opens}
  \ALIUSPlacedTextContent{90.951}{437.733}{418.345}{ALIUSC1F8135}{\ALIUSFontLatoLight}{12.357}{novel ways to directly probe consciousness during sleep. What is the}
  \ALIUSPlacedTextContent{90.951}{455.013}{405.827}{ALIUSC1F8135}{\ALIUSFontLatoLight}{12.357}{relevance of studying sleep according to you to understand consciousness?}
  \ALIUSPlacedTextContent{90.951}{483.928}{418.299}{ALIUSC000000}{\ALIUSFontCormorantRegular}{13.920}{I think there are (at least) three critical points to be made here, and then a}
  \ALIUSPlacedTextContent{90.951}{503.608}{418.303}{ALIUSC000000}{\ALIUSFontCormorantRegular}{13.920}{point specifically about lucid dreaming as methodology. First, sleep is the}
  \ALIUSPlacedTextContent{90.951}{523.049}{418.342}{ALIUSC000000}{\ALIUSFontCormorantRegular}{13.920}{only naturally occurring physiological state characterized by a global loss and}
  \ALIUSPlacedTextContent{90.951}{542.489}{418.320}{ALIUSC000000}{\ALIUSFontCormorantRegular}{13.920}{recovery of consciousness. Thus, sleep turns out to be a useful state for}
  \ALIUSPlacedTextContent{90.951}{561.928}{421.533}{ALIUSC000000}{\ALIUSFontCormorantRegular}{13.920}{building and testing theories of consciousness. Specifically, any}
  \ALIUSPlacedTextContent{90.951}{581.609}{418.218}{ALIUSC000000}{\ALIUSFontCormorantRegular}{13.920}{theory of consciousness needs to both account for and be consistent with the}
  \ALIUSPlacedTextContent{90.951}{601.048}{418.249}{ALIUSC000000}{\ALIUSFontCormorantRegular}{13.920}{neurophysiology of sleep states associated with consciousness. If the}
  \ALIUSPlacedTextContent{90.951}{620.489}{418.245}{ALIUSC000000}{\ALIUSFontCormorantRegular}{13.920}{predictions of a neurobiological theory are inconsistent with the observation}
  \ALIUSPlacedTextContent{90.951}{639.928}{418.257}{ALIUSC000000}{\ALIUSFontCormorantRegular}{13.920}{that consciousness occurs in particular brain states during sleep, then the}
  \ALIUSPlacedTextContent{90.951}{659.609}{418.291}{ALIUSC000000}{\ALIUSFontCormorantRegular}{13.920}{theory needs to be revised or discarded. Having the opportunity to study a}
  \ALIUSPlacedTextContent{90.951}{679.048}{421.533}{ALIUSC000000}{\ALIUSFontCormorantRegular}{13.920}{state with a different neurophysiological milieu seems a great boon to}
  \ALIUSPlacedTextContent{90.951}{698.489}{418.307}{ALIUSC000000}{\ALIUSFontCormorantRegular}{13.920}{scientific research in this area. Being able to study this alternate state puts us}
  \ALIUSPlacedTextContent{90.951}{717.929}{418.193}{ALIUSC000000}{\ALIUSFontCormorantRegular}{13.920}{in a stronger position to build better, more generalizable theories. Second,}
  \ALIUSPlacedTextContent{90.951}{737.369}{421.533}{ALIUSC000000}{\ALIUSFontCormorantRegular}{13.920}{unlike other states, such as anesthesia-induced states of unconsciousness,}
  \ALIUSPlacedTextContent{90.951}{793.143}{122.397}{ALIUSC767171}{\ALIUSFontCormorantLight}{12.000}{ALIUS Bulletin n°4 (2020)}
  \ALIUSPlacedTextContent{387.435}{793.143}{120.728}{ALIUSC767171}{\ALIUSFontCormorantLight}{12.000}{aliusresearch.org/bulletin}
\end{tikzpicture}
\null
\newpage
% Page 10
\thispagestyle{empty}
\noindent\begin{tikzpicture}[remember picture,overlay,x=1bp,y=1bp,shift={(current page.north west)}]
  \ALIUSPlacedTextContent{90.951}{36.329}{322.226}{ALIUSC767171}{\ALIUSFontCormorantRegular}{13.920}{Benjamin Baird – Internally generated conscious activity}
  \ALIUSPlacedTextContent{495.111}{36.329}{14.149}{ALIUSC000000}{\ALIUSFontCormorantRegular}{13.920}{33}
  \ALIUSPlacedTextContent{125.898}{87.307}{14.160}{ALIUSC595959}{\ALIUSFontCormorantLight}{48.000}{\ALIUSPullQuoteOpen}
  \ALIUSPlacedTextContent{143.898}{87.307}{307.204}{ALIUSC000000}{\ALIUSFontLatoLight}{15.120}{What lucid dreaming essentially gives us is}
  \ALIUSPlacedTextContent{143.920}{105.307}{307.160}{ALIUSC000000}{\ALIUSFontLatoLight}{15.120}{experimental control over the dream state in a}
  \ALIUSPlacedTextContent{455.102}{123.307}{14.160}{ALIUSC595959}{\ALIUSFontCormorantLight}{48.000}{\ALIUSPullQuoteClose}
  \ALIUSPlacedTextContent{182.213}{123.307}{230.574}{ALIUSC000000}{\ALIUSFontLatoLight}{15.120}{way that was previously impossible.}
  \ALIUSPlacedTextContent{90.951}{165.208}{418.366}{ALIUSC000000}{\ALIUSFontCormorantRegular}{13.920}{sleep affords us the unique opportunity to contrast consciousness with}
  \ALIUSPlacedTextContent{90.951}{184.889}{418.274}{ALIUSC000000}{\ALIUSFontCormorantRegular}{13.920}{unconsciousness within the same overall vigilance state. There are likely}
  \ALIUSPlacedTextContent{90.951}{204.329}{418.143}{ALIUSC000000}{\ALIUSFontCormorantRegular}{13.920}{many changes in the brain that take place during the shift from wakefulness}
  \ALIUSPlacedTextContent{90.951}{223.768}{418.246}{ALIUSC000000}{\ALIUSFontCormorantRegular}{13.920}{to sleep or anesthesia to post-recovery wakefulness that have nothing to do}
  \ALIUSPlacedTextContent{90.951}{243.209}{108.094}{ALIUSC000000}{\ALIUSFontCormorantRegular}{13.920}{with consciousness}
  \ALIUSPlacedTextContent{200.933}{243.209}{29.406}{ALIUSC000000}{\ALIUSFontCormorantItalic}{13.920}{per se}
  \ALIUSPlacedTextContent{230.368}{243.209}{278.867}{ALIUSC000000}{\ALIUSFontCormorantRegular}{13.920}{. Thus, within-state contrasts in NREM or REM}
  \ALIUSPlacedTextContent{90.951}{262.888}{418.369}{ALIUSC000000}{\ALIUSFontCormorantRegular}{13.920}{sleep allow us to study specific changes in brain activity associated with the}
  \ALIUSPlacedTextContent{90.951}{282.328}{418.249}{ALIUSC000000}{\ALIUSFontCormorantRegular}{13.920}{change in consciousness and to avoid confounding differences in brain}
  \ALIUSPlacedTextContent{90.951}{301.768}{418.179}{ALIUSC000000}{\ALIUSFontCormorantRegular}{13.920}{activity between vigilance states that are not related to consciousness. Third,}
  \ALIUSPlacedTextContent{90.951}{321.208}{418.235}{ALIUSC000000}{\ALIUSFontCormorantRegular}{13.920}{conscious states during sleep are largely decoupled from the external}
  \ALIUSPlacedTextContent{90.951}{340.888}{418.405}{ALIUSC000000}{\ALIUSFontCormorantRegular}{13.920}{environment, so this fortuitously allows us to study the neural substrate of}
  \ALIUSPlacedTextContent{90.951}{360.328}{418.201}{ALIUSC000000}{\ALIUSFontCormorantRegular}{13.920}{consciousness while avoiding another critical confound: namely, the neural}
  \ALIUSPlacedTextContent{90.951}{379.768}{418.090}{ALIUSC000000}{\ALIUSFontCormorantRegular}{13.920}{activity associated with stimulus processing but not related to the generation}
  \ALIUSPlacedTextContent{90.951}{399.208}{125.208}{ALIUSC000000}{\ALIUSFontCormorantRegular}{13.920}{of conscious percepts.}
  \ALIUSPlacedTextContent{90.951}{430.648}{418.172}{ALIUSC000000}{\ALIUSFontCormorantRegular}{13.920}{Finally, in line with the high value of sleep research for consciousness, lucid}
  \ALIUSPlacedTextContent{90.951}{450.328}{418.231}{ALIUSC000000}{\ALIUSFontCormorantRegular}{13.920}{dreaming opens up a new way for us to directly study conscious content}
  \ALIUSPlacedTextContent{90.951}{469.768}{418.217}{ALIUSC000000}{\ALIUSFontCormorantRegular}{13.920}{during ongoing sleep dreaming. I don’t think this fact has yet been}
  \ALIUSPlacedTextContent{90.951}{489.208}{418.196}{ALIUSC000000}{\ALIUSFontCormorantRegular}{13.920}{appreciated or widely recognized. What lucid dreaming essentially gives us is}
  \ALIUSPlacedTextContent{90.951}{508.648}{418.243}{ALIUSC000000}{\ALIUSFontCormorantRegular}{13.920}{experimental control over the dream state in a way that was previously}
  \ALIUSPlacedTextContent{90.951}{528.328}{418.184}{ALIUSC000000}{\ALIUSFontCormorantRegular}{13.920}{impossible. As we have shown in our recent studies, lucid dreamers can}
  \ALIUSPlacedTextContent{90.951}{547.768}{418.249}{ALIUSC000000}{\ALIUSFontCormorantRegular}{13.920}{perform specific tasks during dreams or invoke specific types of conscious}
  \ALIUSPlacedTextContent{90.951}{567.208}{418.304}{ALIUSC000000}{\ALIUSFontCormorantRegular}{13.920}{content (e.g., LaBerge, Baird \& Zimbardo, 2018). Lucid dreaming is thus an}
  \ALIUSPlacedTextContent{90.951}{586.648}{418.322}{ALIUSC000000}{\ALIUSFontCormorantRegular}{13.920}{invaluable methodological tool within the cognitive neuroscience to study}
  \ALIUSPlacedTextContent{90.951}{606.328}{237.739}{ALIUSC000000}{\ALIUSFontCormorantRegular}{13.920}{dreaming and by extension consciousness.}
  \ALIUSPlacedTextContent{90.951}{655.173}{418.323}{ALIUSC1F8135}{\ALIUSFontLatoLight}{12.357}{On a final note, mind-wandering has been reported for accounting about 50\%}
  \ALIUSPlacedTextContent{90.951}{672.453}{418.302}{ALIUSC1F8135}{\ALIUSFontLatoLight}{12.357}{of our awake life (Killingsworth \& Gilbert, 2010) and sleep represents 30\% of}
  \ALIUSPlacedTextContent{90.951}{689.493}{418.397}{ALIUSC1F8135}{\ALIUSFontLatoLight}{12.357}{our lifetime (\ALIUSCitationLink{aliusref-baird-koroma-iber-ancoli-israel-2007}{Iber et al., 2007}). Your research topics cover thus the majority of}
  \ALIUSPlacedTextContent{90.951}{706.773}{421.795}{ALIUSC1F8135}{\ALIUSFontLatoLight}{12.357}{our life. How does research on these topics change our view of mental life}
  \ALIUSPlacedTextContent{90.951}{724.053}{166.366}{ALIUSC1F8135}{\ALIUSFontLatoLight}{12.357}{and the potential of the mind?}
  \ALIUSPlacedTextContent{90.951}{793.143}{122.397}{ALIUSC767171}{\ALIUSFontCormorantLight}{12.000}{ALIUS Bulletin n°4 (2020)}
  \ALIUSPlacedTextContent{387.435}{793.143}{120.728}{ALIUSC767171}{\ALIUSFontCormorantLight}{12.000}{aliusresearch.org/bulletin}
\end{tikzpicture}
\null
\newpage
% Page 11
\thispagestyle{empty}
\noindent\begin{tikzpicture}[remember picture,overlay,x=1bp,y=1bp,shift={(current page.north west)}]
  \ALIUSPlacedTextContent{90.951}{36.329}{322.226}{ALIUSC767171}{\ALIUSFontCormorantRegular}{13.920}{Benjamin Baird – Internally generated conscious activity}
  \ALIUSPlacedTextContent{494.151}{36.329}{15.017}{ALIUSC000000}{\ALIUSFontCormorantRegular}{13.920}{34}
  \ALIUSPlacedTextContent{90.951}{73.048}{418.240}{ALIUSC000000}{\ALIUSFontCormorantRegular}{13.920}{One of the things that I find most fascinating about self-generated states of}
  \ALIUSPlacedTextContent{90.951}{92.488}{418.257}{ALIUSC000000}{\ALIUSFontCormorantRegular}{13.920}{consciousness is that they illustrate the profound extent to which our}
  \ALIUSPlacedTextContent{90.951}{111.928}{418.275}{ALIUSC000000}{\ALIUSFontCormorantRegular}{13.920}{ongoing conscious experiences are generated internally. This is illustrated}
  \ALIUSPlacedTextContent{90.951}{131.609}{418.246}{ALIUSC000000}{\ALIUSFontCormorantRegular}{13.920}{most strikingly in the case of dreams, where, as I mentioned earlier, we can}
  \ALIUSPlacedTextContent{90.951}{151.049}{418.323}{ALIUSC000000}{\ALIUSFontCormorantRegular}{13.920}{experience a multimodal 3D virtual reality that can seem just as real and}
  \ALIUSPlacedTextContent{90.951}{170.488}{418.239}{ALIUSC000000}{\ALIUSFontCormorantRegular}{13.920}{experientially vivid as waking life, but we know is generated completely}
  \ALIUSPlacedTextContent{90.951}{189.928}{418.282}{ALIUSC000000}{\ALIUSFontCormorantRegular}{13.920}{independently of the external environment. For me personally, this has}
  \ALIUSPlacedTextContent{90.951}{209.609}{414.974}{ALIUSC000000}{\ALIUSFontCormorantRegular}{13.920}{profoundly altered my view of perception. I spent most of my life as a naïve-}
  \ALIUSPlacedTextContent{90.951}{229.049}{418.251}{ALIUSC000000}{\ALIUSFontCormorantRegular}{13.920}{realist, thinking that I was looking out through the windows of my eyes and}
  \ALIUSPlacedTextContent{90.951}{248.489}{418.259}{ALIUSC000000}{\ALIUSFontCormorantRegular}{13.920}{seeing directly physical objects in front of me. It turns out that is not how}
  \ALIUSPlacedTextContent{90.951}{267.928}{418.250}{ALIUSC000000}{\ALIUSFontCormorantRegular}{13.920}{perception works. Instead, at all times what we experience is quite literally a}
  \ALIUSPlacedTextContent{90.951}{287.608}{418.210}{ALIUSC000000}{\ALIUSFontCormorantRegular}{13.920}{“virtual reality” generated by our brain. I regard this as one of the most}
  \ALIUSPlacedTextContent{90.951}{307.048}{418.326}{ALIUSC000000}{\ALIUSFontCormorantRegular}{13.920}{interesting things I have ever learned. So overall, one view that emerges from}
  \ALIUSPlacedTextContent{90.951}{326.488}{418.193}{ALIUSC000000}{\ALIUSFontCormorantRegular}{13.920}{studying self-generated states, in particular dreaming, is that our experience}
  \ALIUSPlacedTextContent{90.951}{345.928}{418.303}{ALIUSC000000}{\ALIUSFontCormorantRegular}{13.920}{of the world is at all times endogenously generated by the ongoing dynamics}
  \ALIUSPlacedTextContent{90.951}{365.608}{418.290}{ALIUSC000000}{\ALIUSFontCormorantRegular}{13.920}{of the brain, but is merely shaped by environmental input some of the time.}
  \ALIUSPlacedTextContent{90.951}{385.048}{418.221}{ALIUSC000000}{\ALIUSFontCormorantRegular}{13.920}{This point has been made several times by other researchers, but I emphasize}
  \ALIUSPlacedTextContent{90.951}{404.488}{418.235}{ALIUSC000000}{\ALIUSFontCormorantRegular}{13.920}{it again here because I think it does offer a profound shift in how we think}
  \ALIUSPlacedTextContent{90.951}{423.928}{418.238}{ALIUSC000000}{\ALIUSFontCormorantRegular}{13.920}{about our ongoing conscious mental activities and the nature of our}
  \ALIUSPlacedTextContent{90.951}{443.608}{198.525}{ALIUSC000000}{\ALIUSFontCormorantRegular}{13.920}{perceptual experience of the world.}
  \ALIUSPlacedTextContent{90.951}{475.048}{418.275}{ALIUSC000000}{\ALIUSFontCormorantRegular}{13.920}{In terms of the potentials of the mind, that is a big question. I will just}
  \ALIUSPlacedTextContent{90.951}{494.488}{418.079}{ALIUSC000000}{\ALIUSFontCormorantRegular}{13.920}{mention one here. One of the ways of exploring this that has most fascinated}
  \ALIUSPlacedTextContent{90.951}{513.928}{418.257}{ALIUSC000000}{\ALIUSFontCormorantRegular}{13.920}{me is examining the following question: By extending self-aware}
  \ALIUSPlacedTextContent{90.951}{533.369}{418.271}{ALIUSC000000}{\ALIUSFontCormorantRegular}{13.920}{consciousness into dream and sleep states, could we tap into potentials of our}
  \ALIUSPlacedTextContent{90.951}{553.049}{418.242}{ALIUSC000000}{\ALIUSFontCormorantRegular}{13.920}{minds that we haven’t yet explored? This is yet another one of the reasons I}
  \ALIUSPlacedTextContent{90.951}{572.489}{418.123}{ALIUSC000000}{\ALIUSFontCormorantRegular}{13.920}{have been fascinated by lucid dreaming and drawn to do scientific research}
  \ALIUSPlacedTextContent{90.951}{591.928}{418.254}{ALIUSC000000}{\ALIUSFontCormorantRegular}{13.920}{on the topic. The question is, by being able to bring our “wake-like” cognitive}
  \ALIUSPlacedTextContent{90.951}{611.369}{418.345}{ALIUSC000000}{\ALIUSFontCormorantRegular}{13.920}{faculties into the REM sleep dreaming state, could we explore a part of our}
  \ALIUSPlacedTextContent{90.951}{631.048}{418.272}{ALIUSC000000}{\ALIUSFontCormorantRegular}{13.920}{mind-brain “state-space” that we haven’t previously had access to? And could}
  \ALIUSPlacedTextContent{90.951}{650.489}{418.284}{ALIUSC000000}{\ALIUSFontCormorantRegular}{13.920}{this state have unique uses, for example for things like problem-solving, art}
  \ALIUSPlacedTextContent{90.951}{669.928}{418.317}{ALIUSC000000}{\ALIUSFontCormorantRegular}{13.920}{or creativity? To briefly mention an example, I met a professional composer}
  \ALIUSPlacedTextContent{90.951}{689.369}{418.335}{ALIUSC000000}{\ALIUSFontCormorantRegular}{13.920}{who in his lucid dreams would find a radio and turn it on and hear}
  \ALIUSPlacedTextContent{90.951}{709.048}{418.226}{ALIUSC000000}{\ALIUSFontCormorantRegular}{13.920}{symphonies being played. He would then wake up and transcribe the music}
  \ALIUSPlacedTextContent{90.951}{728.489}{418.290}{ALIUSC000000}{\ALIUSFontCormorantRegular}{13.920}{that he heard in the dream into musical notation. In this way he was able to}
  \ALIUSPlacedTextContent{90.951}{747.929}{418.210}{ALIUSC000000}{\ALIUSFontCormorantRegular}{13.920}{consciously use the REM sleep state for creative inspiration. I’ve met visual}
  \ALIUSPlacedTextContent{90.951}{793.143}{122.397}{ALIUSC767171}{\ALIUSFontCormorantLight}{12.000}{ALIUS Bulletin n°4 (2020)}
  \ALIUSPlacedTextContent{387.435}{793.143}{120.728}{ALIUSC767171}{\ALIUSFontCormorantLight}{12.000}{aliusresearch.org/bulletin}
\end{tikzpicture}
\null
\newpage
% Page 12
\thispagestyle{empty}
\noindent\begin{tikzpicture}[remember picture,overlay,x=1bp,y=1bp,shift={(current page.north west)}]
  \fill[ALIUSC0000FF] (305.511bp,-506.391bp) rectangle ++(200.400bp,-0.720bp);
  \fill[ALIUSC0000FF] (126.951bp,-585.351bp) rectangle ++(206.640bp,-0.720bp);
  \fill[ALIUSC0000FF] (174.711bp,-654.231bp) rectangle ++(218.640bp,-0.720bp);
  \fill[ALIUSC0000FF] (126.951bp,-707.511bp) rectangle ++(232.080bp,-0.720bp);
  \ALIUSPlacedTextContent{90.951}{36.329}{322.226}{ALIUSC767171}{\ALIUSFontCormorantRegular}{13.920}{Benjamin Baird – Internally generated conscious activity}
  \ALIUSPlacedTextContent{494.871}{36.329}{14.429}{ALIUSC000000}{\ALIUSFontCormorantRegular}{13.920}{35}
  \ALIUSPlacedTextContent{90.951}{73.048}{418.226}{ALIUSC000000}{\ALIUSFontCormorantRegular}{13.920}{artists that have done the same type of thing, for instance walking into an}
  \ALIUSPlacedTextContent{90.951}{92.488}{418.188}{ALIUSC000000}{\ALIUSFontCormorantRegular}{13.920}{“art gallery” in their lucid dream and seeing paintings on the walls displayed,}
  \ALIUSPlacedTextContent{90.951}{111.928}{418.265}{ALIUSC000000}{\ALIUSFontCormorantRegular}{13.920}{and then waking up and painting the paintings that they saw in the dream}
  \ALIUSPlacedTextContent{90.951}{131.609}{418.093}{ALIUSC000000}{\ALIUSFontCormorantRegular}{13.920}{gallery. These are anecdotes and this needs to be researched. But it does raise}
  \ALIUSPlacedTextContent{90.951}{151.049}{418.259}{ALIUSC000000}{\ALIUSFontCormorantRegular}{13.920}{the intriguing possibility that there may be untapped potentials of our minds}
  \ALIUSPlacedTextContent{90.951}{170.488}{418.342}{ALIUSC000000}{\ALIUSFontCormorantRegular}{13.920}{that we may be able to access by extending greater awareness into the REM}
  \ALIUSPlacedTextContent{90.951}{189.928}{418.224}{ALIUSC000000}{\ALIUSFontCormorantRegular}{13.920}{sleep state. Overall, I don’t think we’ve really considered yet what the}
  \ALIUSPlacedTextContent{90.951}{209.609}{418.150}{ALIUSC000000}{\ALIUSFontCormorantRegular}{13.920}{potentials could be of making the dreaming state of the brain accessible to}
  \ALIUSPlacedTextContent{90.951}{229.049}{245.774}{ALIUSC000000}{\ALIUSFontCormorantRegular}{13.920}{our unique form of self-aware consciousness.}
  \ALIUSPlacedTextContent{120.439}{294.907}{14.160}{ALIUSC595959}{\ALIUSFontCormorantLight}{48.000}{\ALIUSPullQuoteOpen}
  \ALIUSPlacedTextContent{138.491}{294.907}{318.018}{ALIUSC000000}{\ALIUSFontLatoLight}{15.120}{By extending self-aware consciousness into}
  \ALIUSPlacedTextContent{138.489}{312.907}{318.021}{ALIUSC000000}{\ALIUSFontLatoLight}{15.120}{dream and sleep states, could we tap into}
  \ALIUSPlacedTextContent{138.439}{330.907}{318.122}{ALIUSC000000}{\ALIUSFontLatoLight}{15.120}{potentials of our minds that we haven’t yet}
  \ALIUSPlacedTextContent{266.034}{348.907}{62.932}{ALIUSC000000}{\ALIUSFontLatoLight}{15.120}{explored?}
  \ALIUSPlacedTextContent{460.561}{348.907}{14.160}{ALIUSC595959}{\ALIUSFontCormorantLight}{48.000}{\ALIUSPullQuoteClose}
  \ALIUSPlacedTextContent{264.617}{416.332}{71.454}{ALIUSC000000}{\ALIUSFontLatoLight}{13.920}{References}
  \ALIUSPlacedTextContent{90.951}{446.416}{421.075}{ALIUSC000000}{\ALIUSFontCormorantRegular}{12.960}{\ALIUSRefAnchor{aliusref-baird-koroma-andrillon-windt-2019}Andrillon, T., Windt, J., Silk, T., Drummond, S. P. A., Bellgrove, M. A., \&}
  \ALIUSPlacedTextContent{126.951}{462.016}{385.075}{ALIUSC000000}{\ALIUSFontCormorantRegular}{12.960}{Tsuchiya, N. (2019). Does the Mind Wander When the Brain Takes a Break?}
  \ALIUSPlacedTextContent{126.951}{477.856}{385.075}{ALIUSC000000}{\ALIUSFontCormorantRegular}{12.960}{Local Sleep in Wakefulness, Attentional Lapses and Mind-Wandering.}
  \ALIUSPlacedTextContent{126.951}{493.456}{123.949}{ALIUSC000000}{\ALIUSFontCormorantItalic}{12.960}{Frontiers in Neuroscience}
  \ALIUSPlacedTextContent{250.919}{493.456}{5.919}{ALIUSC000000}{\ALIUSFontCormorantRegular}{12.960}{,}
  \ALIUSPlacedTextContent{260.462}{493.456}{8.215}{ALIUSC000000}{\ALIUSFontCormorantItalic}{12.960}{13}
  \ALIUSPlacedTextContent{268.691}{493.456}{33.102}{ALIUSC000000}{\ALIUSFontCormorantRegular}{12.960}{, 949.}
  \ALIUSPlacedTextContent{305.405}{493.456}{200.560}{ALIUSC0000FF}{\ALIUSFontCormorantRegular}{12.960}{https://doi.org/10.3389/fnins.2019.00949}
  \ALIUSPlacedTextContent{90.951}{525.136}{421.075}{ALIUSC000000}{\ALIUSFontCormorantRegular}{12.960}{\ALIUSRefAnchor{aliusref-baird-koroma-baird-smallwood-2012}Baird, B., Smallwood, J., Mrazek, M. D., Kam, J. W. Y., Franklin, M. S., \& Schooler,}
  \ALIUSPlacedTextContent{126.951}{540.736}{385.075}{ALIUSC000000}{\ALIUSFontCormorantRegular}{12.960}{J. W. (2012). Inspired by Distraction: Mind Wandering Facilitates Creative}
  \ALIUSPlacedTextContent{126.951}{556.576}{61.379}{ALIUSC000000}{\ALIUSFontCormorantRegular}{12.960}{Incubation.}
  \ALIUSPlacedTextContent{188.309}{556.576}{97.472}{ALIUSC000000}{\ALIUSFontCormorantItalic}{12.960}{~Psychological Science}
  \ALIUSPlacedTextContent{285.767}{556.576}{5.919}{ALIUSC000000}{\ALIUSFontCormorantRegular}{12.960}{,}
  \ALIUSPlacedTextContent{291.695}{556.576}{9.151}{ALIUSC000000}{\ALIUSFontCormorantItalic}{12.960}{~23}
  \ALIUSPlacedTextContent{300.860}{556.576}{84.165}{ALIUSC000000}{\ALIUSFontCormorantRegular}{12.960}{(10), 1117–1122.}
  \ALIUSPlacedTextContent{126.951}{572.416}{206.599}{ALIUSC0000FF}{\ALIUSFontCormorantRegular}{12.960}{https://doi.org/10.1177/0956797612446024}
  \ALIUSPlacedTextContent{90.951}{594.016}{421.075}{ALIUSC000000}{\ALIUSFontCormorantRegular}{12.960}{\ALIUSRefAnchor{aliusref-baird-koroma-baird-castelnovo-2018}Baird, B., Castelnovo, A., Gosseries, O., \& Tononi, G. (2018). Frequent lucid}
  \ALIUSPlacedTextContent{126.951}{609.856}{385.075}{ALIUSC000000}{\ALIUSFontCormorantRegular}{12.960}{dreaming associated with increased functional connectivity between}
  \ALIUSPlacedTextContent{126.951}{625.456}{297.475}{ALIUSC000000}{\ALIUSFontCormorantRegular}{12.960}{frontopolar cortex and temporoparietal association areas.}
  \ALIUSPlacedTextContent{425.309}{625.456}{77.770}{ALIUSC000000}{\ALIUSFontCormorantItalic}{12.960}{~Scientific reports}
  \ALIUSPlacedTextContent{503.064}{625.456}{8.961}{ALIUSC000000}{\ALIUSFontCormorantRegular}{12.960}{,}
  \ALIUSPlacedTextContent{126.951}{641.296}{5.702}{ALIUSC000000}{\ALIUSFontCormorantItalic}{12.960}{8}
  \ALIUSPlacedTextContent{132.671}{641.296}{42.019}{ALIUSC000000}{\ALIUSFontCormorantRegular}{12.960}{(1), 1-15.}
  \ALIUSPlacedTextContent{174.699}{641.296}{218.511}{ALIUSC0000FF}{\ALIUSFontCormorantRegular}{12.960}{https://doi.org/ 10.1038/s41598-018-36190-w}
  \ALIUSPlacedTextContent{90.951}{662.896}{421.075}{ALIUSC000000}{\ALIUSFontCormorantRegular}{12.960}{\ALIUSRefAnchor{aliusref-baird-koroma-baird-mota-rolim-2019}Baird, B., Mota-Rolim, S. A., \& Dresler, M. (2019). The cognitive neuroscience of}
  \ALIUSPlacedTextContent{126.951}{678.736}{82.102}{ALIUSC000000}{\ALIUSFontCormorantRegular}{12.960}{lucid dreaming.}
  \ALIUSPlacedTextContent{209.017}{678.736}{180.002}{ALIUSC000000}{\ALIUSFontCormorantItalic}{12.960}{~Neuroscience \& Biobehavioral Reviews}
  \ALIUSPlacedTextContent{389.051}{678.736}{5.919}{ALIUSC000000}{\ALIUSFontCormorantRegular}{12.960}{,}
  \ALIUSPlacedTextContent{394.979}{678.736}{15.148}{ALIUSC000000}{\ALIUSFontCormorantItalic}{12.960}{~100}
  \ALIUSPlacedTextContent{410.150}{678.736}{50.130}{ALIUSC000000}{\ALIUSFontCormorantRegular}{12.960}{, 305–323.}
  \ALIUSPlacedTextContent{126.951}{694.576}{232.247}{ALIUSC0000FF}{\ALIUSFontCormorantRegular}{12.960}{https://doi.org/10.1016/j.neubiorev.2019.03.008}
  \ALIUSPlacedTextContent{90.951}{726.016}{421.075}{ALIUSC000000}{\ALIUSFontCormorantRegular}{12.960}{\ALIUSRefAnchor{aliusref-baird-koroma-baird-riedner-2019}Baird, B., Riedner, B. A., Boly, M., Davidson, R. J., \& Tononi, G. (2019). Increased}
  \ALIUSPlacedTextContent{126.951}{741.616}{385.075}{ALIUSC000000}{\ALIUSFontCormorantRegular}{12.960}{lucid dream frequency in long-term meditators but not following}
  \ALIUSPlacedTextContent{90.951}{793.143}{122.397}{ALIUSC767171}{\ALIUSFontCormorantLight}{12.000}{ALIUS Bulletin n°4 (2020)}
  \ALIUSPlacedTextContent{387.435}{793.143}{120.728}{ALIUSC767171}{\ALIUSFontCormorantLight}{12.000}{aliusresearch.org/bulletin}
\end{tikzpicture}
\null
\newpage
% Page 13
\thispagestyle{empty}
\noindent\begin{tikzpicture}[remember picture,overlay,x=1bp,y=1bp,shift={(current page.north west)}]
  \fill[ALIUSC0000FF] (332.151bp,-101.751bp) rectangle ++(173.760bp,-0.720bp);
  \fill[ALIUSC0000FF] (198.951bp,-170.631bp) rectangle ++(255.120bp,-0.720bp);
  \fill[ALIUSC0000FF] (126.951bp,-249.351bp) rectangle ++(198.720bp,-0.720bp);
  \fill[ALIUSC0000FF] (126.951bp,-328.071bp) rectangle ++(197.520bp,-0.720bp);
  \fill[ALIUSC0000FF] (126.951bp,-406.791bp) rectangle ++(225.360bp,-0.720bp);
  \fill[ALIUSC0000FF] (154.071bp,-469.671bp) rectangle ++(190.080bp,-0.720bp);
  \fill[ALIUSC0000FF] (126.951bp,-564.231bp) rectangle ++(167.520bp,-0.720bp);
  \fill[ALIUSC0000FF] (251.751bp,-648.711bp) rectangle ++(200.640bp,-0.720bp);
  \fill[ALIUSC0000FF] (203.031bp,-727.431bp) rectangle ++(205.200bp,-0.720bp);
  \ALIUSPlacedTextContent{90.951}{36.329}{322.226}{ALIUSC767171}{\ALIUSFontCormorantRegular}{13.920}{Benjamin Baird – Internally generated conscious activity}
  \ALIUSPlacedTextContent{493.911}{36.329}{15.213}{ALIUSC000000}{\ALIUSFontCormorantRegular}{13.920}{36}
  \ALIUSPlacedTextContent{126.951}{72.976}{236.840}{ALIUSC000000}{\ALIUSFontCormorantRegular}{12.960}{mindfulness-based stress reduction training.}
  \ALIUSPlacedTextContent{367.285}{72.976}{144.740}{ALIUSC000000}{\ALIUSFontCormorantItalic}{12.960}{Psychology of Consciousness:}
  \ALIUSPlacedTextContent{126.951}{88.816}{141.683}{ALIUSC000000}{\ALIUSFontCormorantItalic}{12.960}{Theory, Research, and Practice}
  \ALIUSPlacedTextContent{268.677}{88.816}{5.919}{ALIUSC000000}{\ALIUSFontCormorantRegular}{12.960}{,}
  \ALIUSPlacedTextContent{274.039}{88.816}{5.443}{ALIUSC000000}{\ALIUSFontCormorantItalic}{12.960}{~6}
  \ALIUSPlacedTextContent{279.499}{88.816}{53.243}{ALIUSC000000}{\ALIUSFontCormorantRegular}{12.960}{(1), 40–54.}
  \ALIUSPlacedTextContent{332.184}{88.816}{176.799}{ALIUSC0000FF}{\ALIUSFontCormorantRegular}{12.960}{https://doi.org/10.1037/cns0000176}
  \ALIUSPlacedTextContent{90.951}{126.016}{421.075}{ALIUSC000000}{\ALIUSFontCormorantRegular}{12.960}{\ALIUSRefAnchor{aliusref-baird-koroma-brown-engler-1980}Brown, D. P., \& Engler, J. (1980). The stages of mindfulness meditation: A}
  \ALIUSPlacedTextContent{126.951}{141.856}{89.335}{ALIUSC000000}{\ALIUSFontCormorantRegular}{12.960}{validation study.}
  \ALIUSPlacedTextContent{218.557}{141.856}{284.497}{ALIUSC000000}{\ALIUSFontCormorantItalic}{12.960}{The Journal of Transpersonal Psychology; Palo Alto, Calif.}
  \ALIUSPlacedTextContent{503.064}{141.856}{8.961}{ALIUSC000000}{\ALIUSFontCormorantRegular}{12.960}{,}
  \ALIUSPlacedTextContent{126.951}{157.696}{8.488}{ALIUSC000000}{\ALIUSFontCormorantItalic}{12.960}{12}
  \ALIUSPlacedTextContent{135.452}{157.696}{62.129}{ALIUSC000000}{\ALIUSFontCormorantRegular}{12.960}{(2), 143–193.}
  \ALIUSPlacedTextContent{198.951}{157.696}{255.076}{ALIUSC0000FF}{\ALIUSFontCormorantRegular}{12.960}{http://atpweb.org/jtparchive/trps-12-80-02-143.pdf}
  \ALIUSPlacedTextContent{90.951}{189.136}{421.075}{ALIUSC000000}{\ALIUSFontCormorantRegular}{12.960}{\ALIUSRefAnchor{aliusref-baird-koroma-cai-mednick-2009}Cai, D. J., Mednick, S. A., Harrison, E. M., Kanady, J. C., \& Mednick, S. C. (2009).}
  \ALIUSPlacedTextContent{126.951}{204.736}{385.075}{ALIUSC000000}{\ALIUSFontCormorantRegular}{12.960}{REM, not incubation, improves creativity by priming associative networks.}
  \ALIUSPlacedTextContent{126.951}{220.576}{260.431}{ALIUSC000000}{\ALIUSFontCormorantItalic}{12.960}{Proceedings of the National Academy of Sciences}
  \ALIUSPlacedTextContent{387.294}{220.576}{5.919}{ALIUSC000000}{\ALIUSFontCormorantRegular}{12.960}{,}
  \ALIUSPlacedTextContent{399.316}{220.576}{14.914}{ALIUSC000000}{\ALIUSFontCormorantItalic}{12.960}{106}
  \ALIUSPlacedTextContent{414.253}{220.576}{97.772}{ALIUSC000000}{\ALIUSFontCormorantRegular}{12.960}{(25), 10130–10134.}
  \ALIUSPlacedTextContent{126.951}{236.416}{198.681}{ALIUSC0000FF}{\ALIUSFontCormorantRegular}{12.960}{https://doi.org/10.1073/pnas.0900271106}
  \ALIUSPlacedTextContent{90.951}{267.856}{421.075}{ALIUSC000000}{\ALIUSFontCormorantRegular}{12.960}{\ALIUSRefAnchor{aliusref-baird-koroma-colzato-ozturk-2012}Colzato, L. S., Ozturk, A., \& Hommel, B. (2012). Meditate to Create: The Impact}
  \ALIUSPlacedTextContent{126.951}{283.456}{385.075}{ALIUSC000000}{\ALIUSFontCormorantRegular}{12.960}{of Focused-Attention and Open-Monitoring Training on Convergent and}
  \ALIUSPlacedTextContent{126.951}{299.296}{106.452}{ALIUSC000000}{\ALIUSFontCormorantRegular}{12.960}{Divergent Thinking.}
  \ALIUSPlacedTextContent{233.444}{299.296}{107.685}{ALIUSC000000}{\ALIUSFontCormorantItalic}{12.960}{~Frontiers in Psychology}
  \ALIUSPlacedTextContent{341.159}{299.296}{5.919}{ALIUSC000000}{\ALIUSFontCormorantRegular}{12.960}{,}
  \ALIUSPlacedTextContent{347.087}{299.296}{4.432}{ALIUSC000000}{\ALIUSFontCormorantItalic}{12.960}{~3}
  \ALIUSPlacedTextContent{351.533}{299.296}{5.593}{ALIUSC000000}{\ALIUSFontCormorantRegular}{12.960}{.}
  \ALIUSPlacedTextContent{126.951}{315.136}{197.452}{ALIUSC0000FF}{\ALIUSFontCormorantRegular}{12.960}{https://doi.org/10.3389/fpsyg.2012.00116}
  \ALIUSPlacedTextContent{90.951}{346.576}{421.075}{ALIUSC000000}{\ALIUSFontCormorantRegular}{12.960}{\ALIUSRefAnchor{aliusref-baird-koroma-dement-kleitman-1957}Dement, W., \& Kleitman, N. (1957). Cyclic variations in EEG during sleep and}
  \ALIUSPlacedTextContent{126.951}{362.176}{385.075}{ALIUSC000000}{\ALIUSFontCormorantRegular}{12.960}{their relation to eye movements, body motility, and dreaming.}
  \ALIUSPlacedTextContent{126.951}{378.016}{111.407}{ALIUSC000000}{\ALIUSFontCormorantItalic}{12.960}{Electroencephalography}
  \ALIUSPlacedTextContent{249.485}{378.016}{20.500}{ALIUSC000000}{\ALIUSFontCormorantItalic}{12.960}{and}
  \ALIUSPlacedTextContent{281.112}{378.016}{39.396}{ALIUSC000000}{\ALIUSFontCormorantItalic}{12.960}{Clinical}
  \ALIUSPlacedTextContent{331.636}{378.016}{76.638}{ALIUSC000000}{\ALIUSFontCormorantItalic}{12.960}{Neurophysiology}
  \ALIUSPlacedTextContent{408.281}{378.016}{5.919}{ALIUSC000000}{\ALIUSFontCormorantRegular}{12.960}{,}
  \ALIUSPlacedTextContent{425.325}{378.016}{5.443}{ALIUSC000000}{\ALIUSFontCormorantItalic}{12.960}{9}
  \ALIUSPlacedTextContent{430.785}{378.016}{19.811}{ALIUSC000000}{\ALIUSFontCormorantRegular}{12.960}{(4),}
  \ALIUSPlacedTextContent{461.727}{378.016}{50.298}{ALIUSC000000}{\ALIUSFontCormorantRegular}{12.960}{673–690.}
  \ALIUSPlacedTextContent{126.951}{393.856}{225.412}{ALIUSC0000FF}{\ALIUSFontCormorantRegular}{12.960}{https://doi.org/10.1016/0013-4694(57)90088-3}
  \ALIUSPlacedTextContent{90.951}{425.296}{421.075}{ALIUSC000000}{\ALIUSFontCormorantRegular}{12.960}{\ALIUSRefAnchor{aliusref-baird-koroma-ding-tang-2014}Ding, X., Tang, Y.-Y., Tang, R., \& Posner, M. I. (2014). Improving creativity}
  \ALIUSPlacedTextContent{126.951}{440.896}{202.611}{ALIUSC000000}{\ALIUSFontCormorantRegular}{12.960}{performance by short-term meditation.}
  \ALIUSPlacedTextContent{329.677}{440.896}{173.397}{ALIUSC000000}{\ALIUSFontCormorantItalic}{12.960}{~Behavioral and Brain Functions: BBF}
  \ALIUSPlacedTextContent{503.064}{440.896}{8.961}{ALIUSC000000}{\ALIUSFontCormorantRegular}{12.960}{,}
  \ALIUSPlacedTextContent{126.951}{456.736}{9.460}{ALIUSC000000}{\ALIUSFontCormorantItalic}{12.960}{10}
  \ALIUSPlacedTextContent{136.427}{456.736}{17.571}{ALIUSC000000}{\ALIUSFontCormorantRegular}{12.960}{, 9.}
  \ALIUSPlacedTextContent{154.003}{456.736}{190.128}{ALIUSC0000FF}{\ALIUSFontCormorantRegular}{12.960}{https://doi.org/10.1186/1744-9081-10-9}
  \ALIUSPlacedTextContent{90.951}{488.176}{421.075}{ALIUSC000000}{\ALIUSFontCormorantRegular}{12.960}{\ALIUSRefAnchor{aliusref-baird-koroma-dresler-wehrle-2012}Dresler, M., Wehrle, R., Spoormaker, V. I., Koch, S. P., Holsboer, F., Steiger,}
  \ALIUSPlacedTextContent{126.951}{504.016}{382.042}{ALIUSC000000}{\ALIUSFontCormorantRegular}{12.960}{A., Obrig., H., Sämann, P. G., \& Czisch, M. (2012). Neural correlates of}
  \ALIUSPlacedTextContent{126.951}{519.616}{382.033}{ALIUSC000000}{\ALIUSFontCormorantRegular}{12.960}{dream lucidity obtained from contrasting lucid versus non-lucid REM sleep:}
  \ALIUSPlacedTextContent{126.951}{535.456}{8.479}{ALIUSC000000}{\ALIUSFontCormorantRegular}{12.960}{a}
  \ALIUSPlacedTextContent{148.965}{535.456}{53.362}{ALIUSC000000}{\ALIUSFontCormorantRegular}{12.960}{combined}
  \ALIUSPlacedTextContent{215.861}{535.456}{58.610}{ALIUSC000000}{\ALIUSFontCormorantRegular}{12.960}{EEG/fMRI}
  \ALIUSPlacedTextContent{288.005}{535.456}{23.525}{ALIUSC000000}{\ALIUSFontCormorantRegular}{12.960}{case}
  \ALIUSPlacedTextContent{325.065}{535.456}{33.097}{ALIUSC000000}{\ALIUSFontCormorantRegular}{12.960}{study.}
  \ALIUSPlacedTextContent{371.716}{535.456}{23.284}{ALIUSC000000}{\ALIUSFontCormorantItalic}{12.960}{Sleep}
  \ALIUSPlacedTextContent{395.011}{535.456}{5.919}{ALIUSC000000}{\ALIUSFontCormorantRegular}{12.960}{,}
  \ALIUSPlacedTextContent{414.465}{535.456}{9.292}{ALIUSC000000}{\ALIUSFontCormorantItalic}{12.960}{35}
  \ALIUSPlacedTextContent{423.773}{535.456}{19.526}{ALIUSC000000}{\ALIUSFontCormorantRegular}{12.960}{(7),}
  \ALIUSPlacedTextContent{456.840}{535.456}{49.114}{ALIUSC000000}{\ALIUSFontCormorantRegular}{12.960}{1017-1020.}
  \ALIUSPlacedTextContent{126.951}{551.296}{170.518}{ALIUSC0000FF}{\ALIUSFontCormorantRegular}{12.960}{https://doi.org/10.5665/sleep.1974}
  \ALIUSPlacedTextContent{90.951}{588.736}{421.075}{ALIUSC000000}{\ALIUSFontCormorantRegular}{12.960}{\ALIUSRefAnchor{aliusref-baird-koroma-dresler-eibl-2014}Dresler, M., Eibl, L., Fischer, C. F. J., Wehrle, R., Spoormaker, V. I., Steiger, A.,}
  \ALIUSPlacedTextContent{126.951}{604.336}{391.159}{ALIUSC000000}{\ALIUSFontCormorantRegular}{12.960}{Czisch, M., \& Pawlowski, M. (2014). Volitional components of}
  \ALIUSPlacedTextContent{126.951}{620.176}{385.075}{ALIUSC000000}{\ALIUSFontCormorantRegular}{12.960}{consciousness vary across wakefulness, dreaming and lucid dreaming.}
  \ALIUSPlacedTextContent{126.951}{635.776}{107.685}{ALIUSC000000}{\ALIUSFontCormorantItalic}{12.960}{Frontiers in Psychology}
  \ALIUSPlacedTextContent{234.666}{635.776}{5.919}{ALIUSC000000}{\ALIUSFontCormorantRegular}{12.960}{,}
  \ALIUSPlacedTextContent{240.594}{635.776}{5.482}{ALIUSC000000}{\ALIUSFontCormorantItalic}{12.960}{~4}
  \ALIUSPlacedTextContent{246.092}{635.776}{5.594}{ALIUSC000000}{\ALIUSFontCormorantRegular}{12.960}{.}
  \ALIUSPlacedTextContent{251.695}{635.776}{200.689}{ALIUSC0000FF}{\ALIUSFontCormorantRegular}{12.960}{https://doi.org/10.3389/fpsyg.2013.00987}
  \ALIUSPlacedTextContent{90.951}{667.456}{421.075}{ALIUSC000000}{\ALIUSFontCormorantRegular}{12.960}{\ALIUSRefAnchor{aliusref-baird-koroma-fox-nijeboer-2013}Fox, K. C. R., Nijeboer, S., Solomonova, E., Domhoff, G. W., \& Christoff, K.}
  \ALIUSPlacedTextContent{126.951}{683.056}{385.075}{ALIUSC000000}{\ALIUSFontCormorantRegular}{12.960}{(2013). Dreaming as mind wandering: Evidence from functional}
  \ALIUSPlacedTextContent{126.951}{698.896}{267.287}{ALIUSC000000}{\ALIUSFontCormorantRegular}{12.960}{neuroimaging and first-person content reports.}
  \ALIUSPlacedTextContent{400.706}{698.896}{111.319}{ALIUSC000000}{\ALIUSFontCormorantItalic}{12.960}{Frontiers in Human}
  \ALIUSPlacedTextContent{126.951}{714.496}{59.603}{ALIUSC000000}{\ALIUSFontCormorantItalic}{12.960}{Neuroscience}
  \ALIUSPlacedTextContent{186.580}{714.496}{5.919}{ALIUSC000000}{\ALIUSFontCormorantRegular}{12.960}{,}
  \ALIUSPlacedTextContent{192.508}{714.496}{5.016}{ALIUSC000000}{\ALIUSFontCormorantItalic}{12.960}{~7}
  \ALIUSPlacedTextContent{197.539}{714.496}{5.593}{ALIUSC000000}{\ALIUSFontCormorantRegular}{12.960}{.}
  \ALIUSPlacedTextContent{203.142}{714.496}{208.193}{ALIUSC0000FF}{\ALIUSFontCormorantRegular}{12.960}{https://doi.org/10.3389/fnhum.2013.00412}
  \ALIUSPlacedTextContent{90.951}{746.176}{421.075}{ALIUSC000000}{\ALIUSFontCormorantRegular}{12.960}{\ALIUSRefAnchor{aliusref-baird-koroma-grecucci-pappaianni-2015}Grecucci, A., Pappaianni, E., Siugzdaite, R., Theuninck, A., \& Job, R. (2015).}
  \ALIUSPlacedTextContent{90.951}{793.143}{122.397}{ALIUSC767171}{\ALIUSFontCormorantLight}{12.000}{ALIUS Bulletin n°4 (2020)}
  \ALIUSPlacedTextContent{387.435}{793.143}{120.728}{ALIUSC767171}{\ALIUSFontCormorantLight}{12.000}{aliusresearch.org/bulletin}
\end{tikzpicture}
\null
\newpage
% Page 14
\thispagestyle{empty}
\noindent\begin{tikzpicture}[remember picture,overlay,x=1bp,y=1bp,shift={(current page.north west)}]
  \fill[ALIUSC0000FF] (126.951bp,-117.351bp) rectangle ++(174.720bp,-0.720bp);
  \fill[ALIUSC0000FF] (126.951bp,-196.071bp) rectangle ++(261.600bp,-0.720bp);
  \fill[ALIUSC0000FF] (126.951bp,-258.951bp) rectangle ++(162.720bp,-0.720bp);
  \fill[ALIUSC0000FF] (315.351bp,-306.231bp) rectangle ++(190.560bp,-0.720bp);
  \fill[ALIUSC0000FF] (126.951bp,-384.951bp) rectangle ++(228.960bp,-0.720bp);
  \fill[ALIUSC0000FF] (126.951bp,-463.671bp) rectangle ++(213.360bp,-0.720bp);
  \fill[ALIUSC0000FF] (126.951bp,-542.391bp) rectangle ++(220.080bp,-0.720bp);
  \fill[ALIUSC0000FF] (126.951bp,-621.111bp) rectangle ++(202.800bp,-0.720bp);
  \fill[ALIUSC0000FF] (126.951bp,-747.111bp) rectangle ++(229.920bp,-0.720bp);
  \ALIUSPlacedTextContent{90.951}{36.329}{322.226}{ALIUSC767171}{\ALIUSFontCormorantRegular}{13.920}{Benjamin Baird – Internally generated conscious activity}
  \ALIUSPlacedTextContent{494.391}{36.329}{14.709}{ALIUSC000000}{\ALIUSFontCormorantRegular}{13.920}{37}
  \ALIUSPlacedTextContent{126.951}{72.976}{385.075}{ALIUSC000000}{\ALIUSFontCormorantRegular}{12.960}{Mindful Emotion Regulation: Exploring the Neurocognitive Mechanisms}
  \ALIUSPlacedTextContent{126.951}{88.576}{38.382}{ALIUSC000000}{\ALIUSFontCormorantRegular}{12.960}{behind}
  \ALIUSPlacedTextContent{178.430}{88.576}{67.913}{ALIUSC000000}{\ALIUSFontCormorantRegular}{12.960}{Mindfulness.}
  \ALIUSPlacedTextContent{259.463}{88.576}{37.942}{ALIUSC000000}{\ALIUSFontCormorantItalic}{12.960}{BioMed}
  \ALIUSPlacedTextContent{310.515}{88.576}{44.165}{ALIUSC000000}{\ALIUSFontCormorantItalic}{12.960}{Research}
  \ALIUSPlacedTextContent{367.777}{88.576}{61.420}{ALIUSC000000}{\ALIUSFontCormorantItalic}{12.960}{International}
  \ALIUSPlacedTextContent{429.233}{88.576}{5.919}{ALIUSC000000}{\ALIUSFontCormorantRegular}{12.960}{,}
  \ALIUSPlacedTextContent{448.254}{88.576}{19.042}{ALIUSC000000}{\ALIUSFontCormorantItalic}{12.960}{2015}
  \ALIUSPlacedTextContent{467.312}{88.576}{5.919}{ALIUSC000000}{\ALIUSFontCormorantRegular}{12.960}{,}
  \ALIUSPlacedTextContent{486.330}{88.576}{25.696}{ALIUSC000000}{\ALIUSFontCormorantRegular}{12.960}{1–9.}
  \ALIUSPlacedTextContent{126.951}{104.416}{174.640}{ALIUSC0000FF}{\ALIUSFontCormorantRegular}{12.960}{https://doi.org/10.1155/2015/670724}
  \ALIUSPlacedTextContent{90.951}{135.856}{323.554}{ALIUSC000000}{\ALIUSFontCormorantRegular}{12.960}{\ALIUSRefAnchor{aliusref-baird-koroma-iber-ancoli-israel-2007}Iber, C., Ancoli-Israel, S., Chesson, A. L., \& Quan, S. F. (2007).}
  \ALIUSPlacedTextContent{414.954}{135.856}{97.071}{ALIUSC000000}{\ALIUSFontCormorantItalic}{12.960}{The AASM manual}
  \ALIUSPlacedTextContent{126.951}{151.696}{385.075}{ALIUSC000000}{\ALIUSFontCormorantItalic}{12.960}{for the scoring of sleep and associated events: Rules, terminology and technical}
  \ALIUSPlacedTextContent{126.951}{167.296}{382.033}{ALIUSC000000}{\ALIUSFontCormorantItalic}{12.960}{specifications (Vol. 1). Westchester, IL: American Academy of Sleep Medicine.}
  \ALIUSPlacedTextContent{126.951}{183.136}{261.469}{ALIUSC0000FF}{\ALIUSFontCormorantRegular}{12.960}{https://aasm.org/clinical-resources/scoring-manual/}
  \ALIUSPlacedTextContent{90.951}{214.576}{421.075}{ALIUSC000000}{\ALIUSFontCormorantRegular}{12.960}{\ALIUSRefAnchor{aliusref-baird-koroma-kahan-laberge-1994}Kahan, T. L., \& LaBerge, S. (1994). Lucid dreaming as metacognition: Implications}
  \ALIUSPlacedTextContent{126.951}{230.416}{17.959}{ALIUSC000000}{\ALIUSFontCormorantRegular}{12.960}{for}
  \ALIUSPlacedTextContent{155.976}{230.416}{49.534}{ALIUSC000000}{\ALIUSFontCormorantRegular}{12.960}{cognitive}
  \ALIUSPlacedTextContent{216.577}{230.416}{41.539}{ALIUSC000000}{\ALIUSFontCormorantRegular}{12.960}{science.}
  \ALIUSPlacedTextContent{269.183}{230.416}{66.983}{ALIUSC000000}{\ALIUSFontCormorantItalic}{12.960}{Consciousness}
  \ALIUSPlacedTextContent{347.231}{230.416}{20.495}{ALIUSC000000}{\ALIUSFontCormorantItalic}{12.960}{and}
  \ALIUSPlacedTextContent{378.791}{230.416}{45.210}{ALIUSC000000}{\ALIUSFontCormorantItalic}{12.960}{Cognition}
  \ALIUSPlacedTextContent{424.025}{230.416}{5.919}{ALIUSC000000}{\ALIUSFontCormorantRegular}{12.960}{,}
  \ALIUSPlacedTextContent{441.003}{230.416}{4.432}{ALIUSC000000}{\ALIUSFontCormorantItalic}{12.960}{3}
  \ALIUSPlacedTextContent{445.449}{230.416}{5.919}{ALIUSC000000}{\ALIUSFontCormorantRegular}{12.960}{,}
  \ALIUSPlacedTextContent{462.432}{230.416}{43.506}{ALIUSC000000}{\ALIUSFontCormorantRegular}{12.960}{246–264.}
  \ALIUSPlacedTextContent{126.951}{246.016}{162.733}{ALIUSC0000FF}{\ALIUSFontCormorantRegular}{12.960}{https://doi.org/10.1111/cogs.12491}
  \ALIUSPlacedTextContent{90.951}{277.456}{421.075}{ALIUSC000000}{\ALIUSFontCormorantRegular}{12.960}{\ALIUSRefAnchor{aliusref-baird-koroma-killingsworth-gilbert-2010}Killingsworth, M. A., \& Gilbert, D. T. (2010). A Wandering Mind Is an Unhappy}
  \ALIUSPlacedTextContent{126.951}{293.296}{33.385}{ALIUSC000000}{\ALIUSFontCormorantRegular}{12.960}{Mind.}
  \ALIUSPlacedTextContent{164.804}{293.296}{33.144}{ALIUSC000000}{\ALIUSFontCormorantItalic}{12.960}{Science}
  \ALIUSPlacedTextContent{197.954}{293.296}{5.919}{ALIUSC000000}{\ALIUSFontCormorantRegular}{12.960}{,}
  \ALIUSPlacedTextContent{208.339}{293.296}{14.567}{ALIUSC000000}{\ALIUSFontCormorantItalic}{12.960}{330}
  \ALIUSPlacedTextContent{222.925}{293.296}{87.856}{ALIUSC000000}{\ALIUSFontCormorantRegular}{12.960}{(6006), 932–932.}
  \ALIUSPlacedTextContent{315.246}{293.296}{190.671}{ALIUSC0000FF}{\ALIUSFontCormorantRegular}{12.960}{https://doi.org/10.1126/science.1192439}
  \ALIUSPlacedTextContent{90.951}{324.736}{421.075}{ALIUSC000000}{\ALIUSFontCormorantRegular}{12.960}{\ALIUSRefAnchor{aliusref-baird-koroma-kircanski-thompson-2018}Kircanski, K., Thompson, R. J., Sorenson, J., Sherdell, L., \& Gotlib, I. H. (2018).}
  \ALIUSPlacedTextContent{126.951}{340.576}{385.075}{ALIUSC000000}{\ALIUSFontCormorantRegular}{12.960}{The everyday dynamics of rumination and worry: Precipitant events and}
  \ALIUSPlacedTextContent{126.951}{356.176}{45.952}{ALIUSC000000}{\ALIUSFontCormorantRegular}{12.960}{affective}
  \ALIUSPlacedTextContent{183.832}{356.176}{73.391}{ALIUSC000000}{\ALIUSFontCormorantRegular}{12.960}{consequences.}
  \ALIUSPlacedTextContent{268.151}{356.176}{48.266}{ALIUSC000000}{\ALIUSFontCormorantItalic}{12.960}{Cognition}
  \ALIUSPlacedTextContent{327.339}{356.176}{20.495}{ALIUSC000000}{\ALIUSFontCormorantItalic}{12.960}{and}
  \ALIUSPlacedTextContent{358.757}{356.176}{38.877}{ALIUSC000000}{\ALIUSFontCormorantItalic}{12.960}{Emotion}
  \ALIUSPlacedTextContent{397.678}{356.176}{5.919}{ALIUSC000000}{\ALIUSFontCormorantRegular}{12.960}{,}
  \ALIUSPlacedTextContent{414.525}{356.176}{9.150}{ALIUSC000000}{\ALIUSFontCormorantItalic}{12.960}{32}
  \ALIUSPlacedTextContent{423.690}{356.176}{19.526}{ALIUSC000000}{\ALIUSFontCormorantRegular}{12.960}{(7),}
  \ALIUSPlacedTextContent{454.139}{356.176}{57.886}{ALIUSC000000}{\ALIUSFontCormorantRegular}{12.960}{1424–1436.}
  \ALIUSPlacedTextContent{126.951}{372.016}{229.046}{ALIUSC0000FF}{\ALIUSFontCormorantRegular}{12.960}{https://doi.org/10.1080/02699931.2017.1278679}
  \ALIUSPlacedTextContent{90.951}{403.456}{421.075}{ALIUSC000000}{\ALIUSFontCormorantRegular}{12.960}{\ALIUSRefAnchor{aliusref-baird-koroma-laberge-baird-2018}LaBerge, S., Baird, B., \& Zimbardo, P. G. (2018). Smooth tracking of visual targets}
  \ALIUSPlacedTextContent{126.951}{419.296}{385.075}{ALIUSC000000}{\ALIUSFontCormorantRegular}{12.960}{distinguishes lucid REM sleep dreaming and waking perception from}
  \ALIUSPlacedTextContent{126.951}{434.896}{66.933}{ALIUSC000000}{\ALIUSFontCormorantRegular}{12.960}{imagination.}
  \ALIUSPlacedTextContent{193.860}{434.896}{110.770}{ALIUSC000000}{\ALIUSFontCormorantItalic}{12.960}{~Nature communications}
  \ALIUSPlacedTextContent{304.618}{434.896}{5.919}{ALIUSC000000}{\ALIUSFontCormorantRegular}{12.960}{,}
  \ALIUSPlacedTextContent{310.546}{434.896}{5.443}{ALIUSC000000}{\ALIUSFontCormorantItalic}{12.960}{~9}
  \ALIUSPlacedTextContent{316.006}{434.896}{38.743}{ALIUSC000000}{\ALIUSFontCormorantRegular}{12.960}{(1), 1-8.}
  \ALIUSPlacedTextContent{126.951}{450.736}{213.331}{ALIUSC0000FF}{\ALIUSFontCormorantRegular}{12.960}{https://doi.org/10.1038/s41467-018-05547-0}
  \ALIUSPlacedTextContent{90.951}{482.176}{421.075}{ALIUSC000000}{\ALIUSFontCormorantRegular}{12.960}{\ALIUSRefAnchor{aliusref-baird-koroma-laberge-lamarca-2018}LaBerge, S., LaMarca, K., \& Baird, B. (2018). Pre-sleep treatment with galantamine}
  \ALIUSPlacedTextContent{126.951}{498.016}{385.075}{ALIUSC000000}{\ALIUSFontCormorantRegular}{12.960}{stimulates lucid dreaming: A double-blind, placebo-controlled, crossover}
  \ALIUSPlacedTextContent{126.951}{513.616}{33.113}{ALIUSC000000}{\ALIUSFontCormorantRegular}{12.960}{study.}
  \ALIUSPlacedTextContent{160.060}{513.616}{46.170}{ALIUSC000000}{\ALIUSFontCormorantItalic}{12.960}{~PLoS One}
  \ALIUSPlacedTextContent{206.236}{513.616}{5.919}{ALIUSC000000}{\ALIUSFontCormorantRegular}{12.960}{,}
  \ALIUSPlacedTextContent{212.164}{513.616}{8.215}{ALIUSC000000}{\ALIUSFontCormorantItalic}{12.960}{~13}
  \ALIUSPlacedTextContent{220.392}{513.616}{70.293}{ALIUSC000000}{\ALIUSFontCormorantRegular}{12.960}{(8), e0201246.}
  \ALIUSPlacedTextContent{126.951}{529.456}{223.038}{ALIUSC0000FF}{\ALIUSFontCormorantRegular}{12.960}{https://doi.org/10.1371/journal.pone.0201246}
  \ALIUSPlacedTextContent{90.951}{560.896}{421.075}{ALIUSC000000}{\ALIUSFontCormorantRegular}{12.960}{\ALIUSRefAnchor{aliusref-baird-koroma-la-berge-nagel-1981}La Berge, S. P., Nagel, L. E., Dement, W. C., \& Zarcone, V. P. (1981). Lucid}
  \ALIUSPlacedTextContent{126.951}{576.736}{385.075}{ALIUSC000000}{\ALIUSFontCormorantRegular}{12.960}{Dreaming Verified by Volitional Communication during Rem Sleep.}
  \ALIUSPlacedTextContent{126.951}{592.336}{127.669}{ALIUSC000000}{\ALIUSFontCormorantItalic}{12.960}{Perceptual and Motor Skills}
  \ALIUSPlacedTextContent{254.634}{592.336}{5.919}{ALIUSC000000}{\ALIUSFontCormorantRegular}{12.960}{,}
  \ALIUSPlacedTextContent{260.562}{592.336}{9.566}{ALIUSC000000}{\ALIUSFontCormorantItalic}{12.960}{~52}
  \ALIUSPlacedTextContent{270.143}{592.336}{78.883}{ALIUSC000000}{\ALIUSFontCormorantRegular}{12.960}{(3), 727–732.}
  \ALIUSPlacedTextContent{126.951}{608.176}{202.789}{ALIUSC0000FF}{\ALIUSFontCormorantRegular}{12.960}{https://doi.org/10.2466/pms.1981.52.3.727}
  \ALIUSPlacedTextContent{90.951}{639.616}{192.101}{ALIUSC000000}{\ALIUSFontCormorantRegular}{12.960}{\ALIUSRefAnchor{aliusref-baird-koroma-laberge-rheingold-1990}LaBerge, S., \& Rheingold, H. (1990).}
  \ALIUSPlacedTextContent{284.612}{639.616}{193.249}{ALIUSC000000}{\ALIUSFontCormorantItalic}{12.960}{Exploring the World of Lucid Dreaming.}
  \ALIUSPlacedTextContent{477.842}{639.616}{46.351}{ALIUSC000000}{\ALIUSFontCormorantRegular}{12.960}{~New}
  \ALIUSPlacedTextContent{126.951}{655.456}{121.524}{ALIUSC000000}{\ALIUSFontCormorantRegular}{12.960}{York: Ballantine Books.}
  \ALIUSPlacedTextContent{90.951}{686.896}{421.075}{ALIUSC000000}{\ALIUSFontCormorantRegular}{12.960}{\ALIUSRefAnchor{aliusref-baird-koroma-levin-nielsen-2009}Levin, R., \& Nielsen, T. (2009). Nightmares, Bad Dreams, and Emotion}
  \ALIUSPlacedTextContent{126.951}{702.496}{385.075}{ALIUSC000000}{\ALIUSFontCormorantRegular}{12.960}{Dysregulation: A Review and New Neurocognitive Model of Dreaming.}
  \ALIUSPlacedTextContent{126.951}{718.336}{39.873}{ALIUSC000000}{\ALIUSFontCormorantItalic}{12.960}{Current}
  \ALIUSPlacedTextContent{185.302}{718.336}{49.740}{ALIUSC000000}{\ALIUSFontCormorantItalic}{12.960}{Directions}
  \ALIUSPlacedTextContent{253.520}{718.336}{12.294}{ALIUSC000000}{\ALIUSFontCormorantItalic}{12.960}{in}
  \ALIUSPlacedTextContent{284.293}{718.336}{64.287}{ALIUSC000000}{\ALIUSFontCormorantItalic}{12.960}{Psychological}
  \ALIUSPlacedTextContent{367.058}{718.336}{33.136}{ALIUSC000000}{\ALIUSFontCormorantItalic}{12.960}{Science}
  \ALIUSPlacedTextContent{400.213}{718.336}{5.919}{ALIUSC000000}{\ALIUSFontCormorantRegular}{12.960}{,}
  \ALIUSPlacedTextContent{424.606}{718.336}{9.485}{ALIUSC000000}{\ALIUSFontCormorantItalic}{12.960}{18}
  \ALIUSPlacedTextContent{434.109}{718.336}{19.176}{ALIUSC000000}{\ALIUSFontCormorantRegular}{12.960}{(2),}
  \ALIUSPlacedTextContent{471.764}{718.336}{40.261}{ALIUSC000000}{\ALIUSFontCormorantRegular}{12.960}{84–88.}
  \ALIUSPlacedTextContent{126.951}{734.176}{229.831}{ALIUSC0000FF}{\ALIUSFontCormorantRegular}{12.960}{https://doi.org/10.1111/j.1467-8721.2009.01614.x}
  \ALIUSPlacedTextContent{90.951}{793.143}{122.397}{ALIUSC767171}{\ALIUSFontCormorantLight}{12.000}{ALIUS Bulletin n°4 (2020)}
  \ALIUSPlacedTextContent{387.435}{793.143}{120.728}{ALIUSC767171}{\ALIUSFontCormorantLight}{12.000}{aliusresearch.org/bulletin}
\end{tikzpicture}
\null
\newpage
% Page 15
\thispagestyle{empty}
\noindent\begin{tikzpicture}[remember picture,overlay,x=1bp,y=1bp,shift={(current page.north west)}]
  \fill[ALIUSC0000FF] (297.831bp,-132.951bp) rectangle ++(208.080bp,-0.720bp);
  \fill[ALIUSC0000FF] (126.951bp,-196.071bp) rectangle ++(204.480bp,-0.720bp);
  \fill[ALIUSC0000FF] (126.951bp,-274.791bp) rectangle ++(155.760bp,-0.720bp);
  \fill[ALIUSC0000FF] (126.951bp,-353.511bp) rectangle ++(176.160bp,-0.720bp);
  \fill[ALIUSC0000FF] (240.711bp,-416.391bp) rectangle ++(161.520bp,-0.720bp);
  \fill[ALIUSC0000FF] (222.951bp,-463.671bp) rectangle ++(174.240bp,-0.720bp);
  \fill[ALIUSC0000FF] (301.911bp,-510.951bp) rectangle ++(177.360bp,-0.720bp);
  \fill[ALIUSC0000FF] (126.951bp,-579.831bp) rectangle ++(161.760bp,-0.720bp);
  \fill[ALIUSC0000FF] (126.951bp,-642.711bp) rectangle ++(207.840bp,-0.720bp);
  \fill[ALIUSC0000FF] (126.951bp,-753.111bp) rectangle ++(237.840bp,-0.720bp);
  \ALIUSPlacedTextContent{90.951}{36.329}{322.226}{ALIUSC767171}{\ALIUSFontCormorantRegular}{13.920}{Benjamin Baird – Internally generated conscious activity}
  \ALIUSPlacedTextContent{493.671}{36.329}{15.549}{ALIUSC000000}{\ALIUSFontCormorantRegular}{13.920}{38}
  \ALIUSPlacedTextContent{90.951}{72.976}{101.158}{ALIUSC000000}{\ALIUSFontCormorantRegular}{12.960}{\ALIUSRefAnchor{aliusref-baird-koroma-malcolm-1958}Malcolm, N. (1958).}
  \ALIUSPlacedTextContent{192.114}{72.976}{88.199}{ALIUSC000000}{\ALIUSFontCormorantItalic}{12.960}{Dreaming Malcolm}
  \ALIUSPlacedTextContent{280.369}{72.976}{8.635}{ALIUSC000000}{\ALIUSFontCormorantRegular}{12.960}{.}
  \ALIUSPlacedTextContent{90.951}{104.416}{387.579}{ALIUSC000000}{\ALIUSFontCormorantRegular}{12.960}{\ALIUSRefAnchor{aliusref-baird-koroma-nielsen-2004}Nielsen, T. A. (2004). Chronobiological features of dream production.}
  \ALIUSPlacedTextContent{482.655}{104.416}{29.370}{ALIUSC000000}{\ALIUSFontCormorantItalic}{12.960}{Sleep}
  \ALIUSPlacedTextContent{126.951}{120.016}{84.849}{ALIUSC000000}{\ALIUSFontCormorantItalic}{12.960}{Medicine Reviews}
  \ALIUSPlacedTextContent{211.786}{120.016}{5.919}{ALIUSC000000}{\ALIUSFontCormorantRegular}{12.960}{,}
  \ALIUSPlacedTextContent{220.663}{120.016}{5.702}{ALIUSC000000}{\ALIUSFontCormorantItalic}{12.960}{8}
  \ALIUSPlacedTextContent{226.383}{120.016}{68.589}{ALIUSC000000}{\ALIUSFontCormorantRegular}{12.960}{(5), 403–424.}
  \ALIUSPlacedTextContent{297.931}{120.016}{207.979}{ALIUSC0000FF}{\ALIUSFontCormorantRegular}{12.960}{https://doi.org/10.1016/j.smrv.2004.06.005}
  \ALIUSPlacedTextContent{90.951}{151.696}{421.075}{ALIUSC000000}{\ALIUSFontCormorantRegular}{12.960}{\ALIUSRefAnchor{aliusref-baird-koroma-schredl-erlacher-2004}Schredl, M., \& Erlacher, D. (2004). Lucid dreaming frequency and personality.}
  \ALIUSPlacedTextContent{126.951}{167.296}{55.112}{ALIUSC000000}{\ALIUSFontCormorantItalic}{12.960}{Personality}
  \ALIUSPlacedTextContent{205.151}{167.296}{20.496}{ALIUSC000000}{\ALIUSFontCormorantItalic}{12.960}{and}
  \ALIUSPlacedTextContent{248.736}{167.296}{51.425}{ALIUSC000000}{\ALIUSFontCormorantItalic}{12.960}{Individual}
  \ALIUSPlacedTextContent{323.249}{167.296}{50.220}{ALIUSC000000}{\ALIUSFontCormorantItalic}{12.960}{Differences}
  \ALIUSPlacedTextContent{373.497}{167.296}{5.919}{ALIUSC000000}{\ALIUSFontCormorantRegular}{12.960}{,}
  \ALIUSPlacedTextContent{402.506}{167.296}{9.461}{ALIUSC000000}{\ALIUSFontCormorantItalic}{12.960}{37}
  \ALIUSPlacedTextContent{411.982}{167.296}{19.526}{ALIUSC000000}{\ALIUSFontCormorantRegular}{12.960}{(7),}
  \ALIUSPlacedTextContent{454.601}{167.296}{51.330}{ALIUSC000000}{\ALIUSFontCormorantRegular}{12.960}{1463–1473.}
  \ALIUSPlacedTextContent{126.951}{183.136}{204.603}{ALIUSC0000FF}{\ALIUSFontCormorantRegular}{12.960}{https://doi.org/10.1016/j.paid.2004.02.003}
  \ALIUSPlacedTextContent{90.951}{214.576}{421.075}{ALIUSC000000}{\ALIUSFontCormorantRegular}{12.960}{\ALIUSRefAnchor{aliusref-baird-koroma-siclari-baird-2017}Siclari, F., Baird, B., Perogamvros, L., Bernardi, G., LaRocque, J. J., Riedner, B.,}
  \ALIUSPlacedTextContent{126.951}{230.416}{385.075}{ALIUSC000000}{\ALIUSFontCormorantRegular}{12.960}{Boly, M., Postle, B. R., \& Tononi, G. (2017). The neural correlates of}
  \ALIUSPlacedTextContent{126.951}{246.016}{53.711}{ALIUSC000000}{\ALIUSFontCormorantRegular}{12.960}{dreaming.}
  \ALIUSPlacedTextContent{217.839}{246.016}{35.886}{ALIUSC000000}{\ALIUSFontCormorantItalic}{12.960}{Nature}
  \ALIUSPlacedTextContent{290.908}{246.016}{59.603}{ALIUSC000000}{\ALIUSFontCormorantItalic}{12.960}{Neuroscience}
  \ALIUSPlacedTextContent{350.549}{246.016}{5.919}{ALIUSC000000}{\ALIUSFontCormorantRegular}{12.960}{,}
  \ALIUSPlacedTextContent{393.651}{246.016}{10.395}{ALIUSC000000}{\ALIUSFontCormorantItalic}{12.960}{20}
  \ALIUSPlacedTextContent{404.064}{246.016}{19.992}{ALIUSC000000}{\ALIUSFontCormorantRegular}{12.960}{(6),}
  \ALIUSPlacedTextContent{461.237}{246.016}{50.789}{ALIUSC000000}{\ALIUSFontCormorantRegular}{12.960}{872–878.}
  \ALIUSPlacedTextContent{126.951}{261.856}{155.718}{ALIUSC0000FF}{\ALIUSFontCormorantRegular}{12.960}{https://doi.org/10.1038/nn.4545}
  \ALIUSPlacedTextContent{90.951}{293.296}{421.075}{ALIUSC000000}{\ALIUSFontCormorantRegular}{12.960}{\ALIUSRefAnchor{aliusref-baird-koroma-smallwood-fishman-2007}Smallwood, J., Fishman, D. J., \& Schooler, J. W. (2007). Counting the cost of an}
  \ALIUSPlacedTextContent{126.951}{309.136}{385.075}{ALIUSC000000}{\ALIUSFontCormorantRegular}{12.960}{absent mind: Mind wandering as an underrecognized influence on}
  \ALIUSPlacedTextContent{126.951}{324.736}{134.968}{ALIUSC000000}{\ALIUSFontCormorantRegular}{12.960}{educational performance.}
  \ALIUSPlacedTextContent{265.111}{324.736}{158.060}{ALIUSC000000}{\ALIUSFontCormorantItalic}{12.960}{Psychonomic Bulletin \& Review}
  \ALIUSPlacedTextContent{423.155}{324.736}{5.919}{ALIUSC000000}{\ALIUSFontCormorantRegular}{12.960}{,}
  \ALIUSPlacedTextContent{432.249}{324.736}{9.265}{ALIUSC000000}{\ALIUSFontCormorantItalic}{12.960}{14}
  \ALIUSPlacedTextContent{441.531}{324.736}{70.495}{ALIUSC000000}{\ALIUSFontCormorantRegular}{12.960}{(2), 230–236.}
  \ALIUSPlacedTextContent{126.951}{340.576}{176.169}{ALIUSC0000FF}{\ALIUSFontCormorantRegular}{12.960}{https://doi.org/10.3758/BF03194057}
  \ALIUSPlacedTextContent{90.951}{372.016}{421.075}{ALIUSC000000}{\ALIUSFontCormorantRegular}{12.960}{\ALIUSRefAnchor{aliusref-baird-koroma-smallwood-fitzgerald-2009}Smallwood, J., Fitzgerald, A., Miles, L. K., \& Phillips, L. H. (2009). Shifting}
  \ALIUSPlacedTextContent{126.951}{387.856}{385.075}{ALIUSC000000}{\ALIUSFontCormorantRegular}{12.960}{moods, wandering minds: Negative moods lead the mind to wander.}
  \ALIUSPlacedTextContent{126.951}{403.456}{38.901}{ALIUSC000000}{\ALIUSFontCormorantItalic}{12.960}{Emotion}
  \ALIUSPlacedTextContent{165.859}{403.456}{5.919}{ALIUSC000000}{\ALIUSFontCormorantRegular}{12.960}{,}
  \ALIUSPlacedTextContent{171.787}{403.456}{5.443}{ALIUSC000000}{\ALIUSFontCormorantItalic}{12.960}{~9}
  \ALIUSPlacedTextContent{177.247}{403.456}{63.432}{ALIUSC000000}{\ALIUSFontCormorantRegular}{12.960}{(2), 271–276.}
  \ALIUSPlacedTextContent{240.698}{403.456}{164.618}{ALIUSC0000FF}{\ALIUSFontCormorantRegular}{12.960}{https://doi.org/10.1037/a0014855}
  \ALIUSPlacedTextContent{90.951}{434.896}{252.707}{ALIUSC000000}{\ALIUSFontCormorantRegular}{12.960}{\ALIUSRefAnchor{aliusref-baird-koroma-tulving-1985}Tulving, E. (1985). Memory and consciousness.}
  \ALIUSPlacedTextContent{346.954}{434.896}{165.071}{ALIUSC000000}{\ALIUSFontCormorantItalic}{12.960}{Canadian Psychology/Psychologie}
  \ALIUSPlacedTextContent{126.951}{450.736}{51.774}{ALIUSC000000}{\ALIUSFontCormorantItalic}{12.960}{canadienne}
  \ALIUSPlacedTextContent{178.755}{450.736}{5.919}{ALIUSC000000}{\ALIUSFontCormorantRegular}{12.960}{,}
  \ALIUSPlacedTextContent{184.683}{450.736}{10.162}{ALIUSC000000}{\ALIUSFontCormorantItalic}{12.960}{~26}
  \ALIUSPlacedTextContent{194.862}{450.736}{25.148}{ALIUSC000000}{\ALIUSFontCormorantRegular}{12.960}{(1), 1.}
  \ALIUSPlacedTextContent{223.016}{451.539}{174.092}{ALIUSC0000FF}{\ALIUSFontCormorantRegular}{12.960}{http://doi.org/10.1037/h0080017}
  \ALIUSPlacedTextContent{90.951}{482.176}{421.075}{ALIUSC000000}{\ALIUSFontCormorantRegular}{12.960}{\ALIUSRefAnchor{aliusref-baird-koroma-wagner-gais-2004}Wagner, U., Gais, S., Haider, H., Verleger, R., \& Born, J. (2004). Sleep inspires}
  \ALIUSPlacedTextContent{126.951}{498.016}{40.479}{ALIUSC000000}{\ALIUSFontCormorantRegular}{12.960}{insight.}
  \ALIUSPlacedTextContent{167.431}{498.016}{32.841}{ALIUSC000000}{\ALIUSFontCormorantItalic}{12.960}{~Nature}
  \ALIUSPlacedTextContent{200.295}{498.016}{5.919}{ALIUSC000000}{\ALIUSFontCormorantRegular}{12.960}{,}
  \ALIUSPlacedTextContent{206.223}{498.016}{15.236}{ALIUSC000000}{\ALIUSFontCormorantItalic}{12.960}{~427}
  \ALIUSPlacedTextContent{221.471}{498.016}{80.426}{ALIUSC000000}{\ALIUSFontCormorantRegular}{12.960}{(6972), 352–355.}
  \ALIUSPlacedTextContent{301.913}{498.016}{177.228}{ALIUSC0000FF}{\ALIUSFontCormorantRegular}{12.960}{https://doi.org/10.1038/nature02223}
  \ALIUSPlacedTextContent{90.951}{535.456}{421.075}{ALIUSC000000}{\ALIUSFontCormorantRegular}{12.960}{\ALIUSRefAnchor{aliusref-baird-koroma-walker-van-der-helm-2009}Walker, M. P., \& van der Helm, E. (2009). Overnight therapy? The role of sleep in}
  \ALIUSPlacedTextContent{126.951}{551.056}{161.237}{ALIUSC000000}{\ALIUSFontCormorantRegular}{12.960}{emotional brain processing.}
  \ALIUSPlacedTextContent{297.380}{551.056}{109.972}{ALIUSC000000}{\ALIUSFontCormorantItalic}{12.960}{Psychological Bulletin}
  \ALIUSPlacedTextContent{407.363}{551.056}{5.919}{ALIUSC000000}{\ALIUSFontCormorantRegular}{12.960}{,}
  \ALIUSPlacedTextContent{422.435}{551.056}{13.074}{ALIUSC000000}{\ALIUSFontCormorantItalic}{12.960}{135}
  \ALIUSPlacedTextContent{435.526}{551.056}{76.500}{ALIUSC000000}{\ALIUSFontCormorantRegular}{12.960}{(5), 731–748.}
  \ALIUSPlacedTextContent{126.951}{566.896}{161.861}{ALIUSC0000FF}{\ALIUSFontCormorantRegular}{12.960}{https://doi.org/10.1037/a0016570}
  \ALIUSPlacedTextContent{90.951}{598.336}{424.117}{ALIUSC000000}{\ALIUSFontCormorantRegular}{12.960}{\ALIUSRefAnchor{aliusref-baird-koroma-windt-2013}Windt, J. M. (2013). Reporting dream experience: Why (not) to be skeptical about}
  \ALIUSPlacedTextContent{126.951}{614.176}{35.096}{ALIUSC000000}{\ALIUSFontCormorantRegular}{12.960}{dream}
  \ALIUSPlacedTextContent{184.280}{614.176}{42.001}{ALIUSC000000}{\ALIUSFontCormorantRegular}{12.960}{reports.}
  \ALIUSPlacedTextContent{248.521}{614.176}{44.783}{ALIUSC000000}{\ALIUSFontCormorantItalic}{12.960}{Frontiers}
  \ALIUSPlacedTextContent{315.538}{614.176}{12.286}{ALIUSC000000}{\ALIUSFontCormorantItalic}{12.960}{in}
  \ALIUSPlacedTextContent{350.057}{614.176}{38.290}{ALIUSC000000}{\ALIUSFontCormorantItalic}{12.960}{Human}
  \ALIUSPlacedTextContent{410.580}{614.176}{59.635}{ALIUSC000000}{\ALIUSFontCormorantItalic}{12.960}{Neuroscience}
  \ALIUSPlacedTextContent{470.209}{614.176}{5.919}{ALIUSC000000}{\ALIUSFontCormorantRegular}{12.960}{,}
  \ALIUSPlacedTextContent{498.359}{614.176}{5.016}{ALIUSC000000}{\ALIUSFontCormorantItalic}{12.960}{7}
  \ALIUSPlacedTextContent{503.390}{614.176}{8.635}{ALIUSC000000}{\ALIUSFontCormorantRegular}{12.960}{.}
  \ALIUSPlacedTextContent{126.951}{629.776}{207.921}{ALIUSC0000FF}{\ALIUSFontCormorantRegular}{12.960}{https://doi.org/10.3389/fnhum.2013.00708}
  \ALIUSPlacedTextContent{90.951}{661.456}{219.421}{ALIUSC000000}{\ALIUSFontCormorantRegular}{12.960}{\ALIUSRefAnchor{aliusref-baird-koroma-windt-metzinger-2007}Windt, J. M., \& Metzinger, T. (2007).}
  \ALIUSPlacedTextContent{315.419}{661.456}{193.564}{ALIUSC000000}{\ALIUSFontCormorantItalic}{12.960}{The philosophy of dreaming and self}
  \ALIUSPlacedTextContent{126.951}{677.056}{385.075}{ALIUSC000000}{\ALIUSFontCormorantItalic}{12.960}{consciousness: what happens to the experiential subject during the dream}
  \ALIUSPlacedTextContent{126.951}{692.896}{28.077}{ALIUSC000000}{\ALIUSFontCormorantItalic}{12.960}{state?.}
  \ALIUSPlacedTextContent{155.043}{692.896}{353.973}{ALIUSC000000}{\ALIUSFontCormorantRegular}{12.960}{~In D. Barrett \& P. McNamara (Eds.), Praeger perspectives. The new}
  \ALIUSPlacedTextContent{126.951}{708.496}{378.979}{ALIUSC000000}{\ALIUSFontCormorantRegular}{12.960}{science of dreaming: Vol. 3. Cultural and theoretical perspectives (p. 193–}
  \ALIUSPlacedTextContent{126.951}{724.336}{26.270}{ALIUSC000000}{\ALIUSFontCormorantRegular}{12.960}{247).}
  \ALIUSPlacedTextContent{178.106}{724.336}{41.399}{ALIUSC000000}{\ALIUSFontCormorantRegular}{12.960}{Praeger}
  \ALIUSPlacedTextContent{244.391}{724.336}{118.947}{ALIUSC000000}{\ALIUSFontCormorantRegular}{12.960}{Publishers/Greenwood}
  \ALIUSPlacedTextContent{388.224}{724.336}{56.972}{ALIUSC000000}{\ALIUSFontCormorantRegular}{12.960}{Publishing}
  \ALIUSPlacedTextContent{470.082}{724.336}{38.901}{ALIUSC000000}{\ALIUSFontCormorantRegular}{12.960}{Group.}
  \ALIUSPlacedTextContent{126.951}{740.176}{237.707}{ALIUSC0000FF}{\ALIUSFontCormorantRegular}{12.960}{https://psycnet.apa.org/record/2007-09897-009}
  \ALIUSPlacedTextContent{90.951}{793.143}{122.397}{ALIUSC767171}{\ALIUSFontCormorantLight}{12.000}{ALIUS Bulletin n°4 (2020)}
  \ALIUSPlacedTextContent{387.435}{793.143}{120.728}{ALIUSC767171}{\ALIUSFontCormorantLight}{12.000}{aliusresearch.org/bulletin}
\end{tikzpicture}
\null
\ifdefined\ALIUSIssueBuild
\clearpage
\expandafter\endinput
\fi
\end{document}
